\documentclass[11pt,oneside]{report}
% ======================================================================
%  THE TRIADIC MEASUREMENT SUBSTRATE --- COMPLETE CORPUS (Volumes 1 & 2)
%  Aaron Switzer, 2026
%  Single-file build. Compile: pdflatex x3 (third pass settles the TOC).
% ======================================================================
% ======================================================================
%  TMS Volume 1 -- shared preamble
%  Compiles with pdfLaTeX both locally and on Overleaf (no exotic fonts).
% ======================================================================
\usepackage[T1]{fontenc}
\usepackage[utf8]{inputenc}
\usepackage{amsmath}
\usepackage{amssymb}
\usepackage{mathptmx}            % Times text + math (widely available)
\usepackage[scaled=0.90]{helvet} % sans for headings
\usepackage{courier}

\usepackage[letterpaper,margin=1in]{geometry}
\usepackage{amsthm}
\usepackage{mathtools}
\usepackage{bm}
\usepackage{microtype}
\usepackage{enumitem}
\usepackage{booktabs}
\usepackage{array}
\usepackage{longtable}
\usepackage{caption}
\usepackage{xcolor}
\usepackage{titlesec}
\usepackage{fancyhdr}
\usepackage[hidelinks]{hyperref}
\usepackage{tocbibind}

% ---- spacing -----------------------------------------------------------
\setlength{\parindent}{1.2em}
\setlength{\parskip}{0.35em}
\linespread{1.04}
\setlist{leftmargin=1.6em,topsep=0.25em,itemsep=0.15em,parsep=0pt}

% ---- theorem environments ---------------------------------------------
\newtheorem{theorem}{Theorem}[section]
\newtheorem{lemma}[theorem]{Lemma}
\newtheorem{corollary}[theorem]{Corollary}
\newtheorem{proposition}[theorem]{Proposition}

\theoremstyle{definition}
\newtheorem{definition}[theorem]{Definition}
\newtheorem{axiom}{Axiom}
\newtheorem{claim}[theorem]{Claim}
\newtheorem{principle}[theorem]{Principle}
\newtheorem{conjecture}[theorem]{Conjecture}
\newtheorem{prediction}[theorem]{Prediction}
\newtheorem{property}[theorem]{Property}
\newtheorem{postulate}[theorem]{Postulate}

\theoremstyle{remark}
\newtheorem{remark}[theorem]{Remark}
\newtheorem*{remark*}{Remark}

% ---- section heading look (unnumbered; author supplies the numbers) ----
\titleformat{\section}[hang]{\normalfont\large\bfseries\sffamily}{}{0pt}{}
\titleformat{\subsection}[hang]{\normalfont\normalsize\bfseries\sffamily}{}{0pt}{}
\titleformat{\subsubsection}[hang]{\normalfont\normalsize\bfseries\itshape}{}{0pt}{}
\titlespacing*{\section}{0pt}{1.4em}{0.5em}
\titlespacing*{\subsection}{0pt}{1.0em}{0.35em}
\titlespacing*{\subsubsection}{0pt}{0.8em}{0.3em}

% chapter (= each paper) heading
\titleformat{\chapter}[display]
  {\normalfont\sffamily\bfseries}
  {}{0pt}{\Huge}
\titlespacing*{\chapter}{0pt}{10pt}{24pt}

% ---- page-break quality: avoid stranded headings and orphaned/widowed lines ----
% (applies to both volumes via the shared preamble)
\widowpenalty=10000        % no single line at the top of a page
\clubpenalty=10000         % no single line at the bottom of a page
\displaywidowpenalty=10000 % same, around displayed equations
\brokenpenalty=4000        % discourage page breaks right after a hyphenated line
% titlesec: forbid a page break immediately after a heading (heading stranded at page foot)
\usepackage{needspace}
\usepackage{etoolbox}
\pretocmd{\section}{\Needspace*{4\baselineskip}}{}{}
\pretocmd{\subsection}{\Needspace*{3\baselineskip}}{}{}

% ---- running heads -----------------------------------------------------
\pagestyle{fancy}
\fancyhf{}
\fancyhead[L]{\small\itshape The Triadic Measurement Substrate}
\fancyhead[R]{\small\itshape\nouppercase{\rightmark}}
\fancyfoot[C]{\small\thepage}
\renewcommand{\headrulewidth}{0.4pt}
\fancypagestyle{plain}{\fancyhf{}\fancyfoot[C]{\small\thepage}\renewcommand{\headrulewidth}{0pt}}

% ---- abstract / keywords / references list ----------------------------
\newenvironment{paperabstract}
  {\begin{center}\begin{minipage}{0.88\textwidth}\small
   \noindent\textbf{Abstract.}\itshape\space}
  {\end{minipage}\end{center}\vspace{0.4em}}

\newcommand{\keywords}[1]{\noindent{\small\textbf{Keywords:} #1}\vspace{0.4em}}

% per-paper APA reference list (hanging indent), kept as authored
\newenvironment{reflist}
  {\par\small\setlength{\parindent}{0pt}%
   \setlength{\leftskip}{1.6em}\setlength{\parindent}{-1.6em}%
   \setlength{\parskip}{0.4em}%
   \raggedright\hyphenpenalty=50\exhyphenpenalty=50\relax
   \sloppy}
  {\par}

% convenience
\providecommand{\tightlist}{\setlength{\itemsep}{0pt}\setlength{\parskip}{0pt}}
\providecommand{\TMS}{\textsc{tms}}
\hyphenation{con-fir-ma-tion sub-strate}
\begin{document}

\thispagestyle{empty}
\begin{titlepage}
\centering
\vspace*{2cm}
{\sffamily\bfseries\Large The Triadic Measurement Substrate\par}
\vspace{0.6cm}
{\large $\bar{B} = 3\bar{\alpha}$\par}
\vspace{1cm}
{\sffamily\large Volume 1\par}
\vspace{1.4cm}
{\Large \bfseries Paper 4: Hierarchies, Self-Modeling Measurement Nodes, Multi-Subsystem Dynamics, and Novelty Exhaustion\par}
\vfill
{\large Aaron Switzer\par}
\vspace{0.4cm}
{\large 2026\par}
\vspace{1.2cm}
{\small\itshape Excerpted from the complete two-volume corpus, \emph{The Triadic Measurement Substrate}. Full corpus and reproducibility package: \url{https://doi.org/10.5281/zenodo.22035422}\par}
\end{titlepage}
\cleardoublepage
\pagenumbering{arabic}

% ---------------- Volume 1 / paper4 ----------------
\chapter*{Paper 4}
\markboth{Paper 4}{Paper 4}
\addcontentsline{toc}{chapter}{Paper 4: Hierarchies, Self-Modeling Measurement Nodes, Multi-Subsystem Dynamics, and Novelty Exhaustion}
\begingroup\centering
{\large\itshape Hierarchies, Self-Modeling Measurement Nodes, Multi-Subsystem Dynamics, and Novelty Exhaustion\par}\vspace{0.8em}
{\normalsize Aaron Switzer\par}
{\small May 2026\par}
\endgroup\vspace{1.0em}

\noindent{\small\textbf{Keywords:} TMS, Subsystem Hierarchy, Recursive Coarse-Graining, Self-Modeling, Measurement Node, Stressor Dynamics, Region Biography, Novelty Exhaustion, Hierarchical Lightspeeds}\par\vspace{0.6em}

\section*{Status of Results}
\addcontentsline{toc}{section}{Status of Results}

This paper distinguishes four categories of claim, each labeled where it
appears in the text.

\textbf{Derived.} Results that follow from the five axioms, the
Principle of Generative Economy, and the structures established in
Papers 1 through 3 by explicit construction or proof. The recursive
coarse-graining operator \ensuremath{G^{(k)}} at arbitrary hierarchical depth as the
recursive application of Paper 3's G under axiom scale-independence
(\S\,2.1); the scale-separation requirement \ensuremath{L_{p}^{(k)}} \ensuremath{\gg} \ensuremath{L_{p}^{(k-1)}} and
\ensuremath{T_{p}^{(k)}} \ensuremath{\gg} \ensuremath{T_{p}^{(k-1)}} as forced by the meta-level coherence
threshold (\S\,2.2); the existence and finiteness of the hierarchy ceiling
\ensuremath{k_{\mathrm{max}}} for any bounded region (\S\,2.3); the Russian-doll floor at \ensuremath{L_{p}} for
downward nesting (\S\,2.4); the structural asymmetry between lossy
coarse-graining and generative downward nesting, establishing per-level
independence of configuration spaces (\S\,2.5.1); the three growth modes
(annexation, imprinting, incorporation) with their respective geometric,
informational, and structural ceilings (\S\,III); self-modeling as a
structural property of every subsystem at k \ensuremath{\geq} 1, with the three
conditions of self-modeling jointly satisfied by the \ensuremath{(C_{D},} \ensuremath{C_{O},} \ensuremath{C_{C})}
trichotomy plus the Imprinting Theorem (\S\,4.2); the identification of
every subsystem at k \ensuremath{\geq} 1 as a measurement node at its level (\S\,4.3);
higher-order self-modeling at k \ensuremath{\geq} 2 including self-modeling-of-self as a
derived consequence of the recursion (\S\,4.6); meta-meta-confirmation as
the same triadic structure applied at the next level, with the
polarization-biased resolution rule inherited unchanged (\S\,5.2); the
coupling-timescale constraint \ensuremath{T_{p}^{(k+1)}} \ensuremath{\geq} \ensuremath{\sqrt{2}} \ensuremath{\cdot} \ensuremath{(L_{p}(k+1)/L_{p}(k))} \ensuremath{\cdot}
\ensuremath{T_{p}^{(k)}} from the requirement of coherent meta-meta-confirmation
(\S\,5.3); three substrate-level stressors (mortality, predation,
starvation) with substrate mechanisms and level-(k+1) signatures
(\S\,5.5--\S\,5.7); the overdetermination of subsystem formation by
Borel-Cantelli inevitability, PGE preference, and flux-balance forcing
(\S\,5.1); inter-region and intra-region diversity as derived consequences
of substrate-level dynamics (\S\,5.9); the finiteness of accessible
configuration space \ensuremath{\Sigma (R)} for any bounded region (\S\,6.1); monotonic decay
of the novelty production rate \ensuremath{\nu (R,} \ensuremath{\tau )} as \ensuremath{\Sigma_{\mathrm{sampled}}} approaches \ensuremath{\Sigma (R)}
(\S\,6.2); the structural exhaustion problem at \ensuremath{\tau} \textgreater{} \ensuremath{\tau *(R),}
with Axioms 3 and 4 jointly forcing reinterpretation as the only
available continuation (\S\,6.6); and the upper bound \ensuremath{c^{(k+1)}} \ensuremath{\leq} \ensuremath{(1/\sqrt{2})} \ensuremath{\cdot}
\ensuremath{c^{(k)}} on the hierarchical lightspeed dispersion as a structural law
of the recursive FCC face-diagonal-to-edge ratio at every hierarchical
level (\S\,7.4 sharpened with reference to Paper 5 \S\,7.2 and Paper 3 \S\,4.5,
holding in all regimes; the tight-coupling case saturates equality) are
derived in this sense.

\textbf{Structural correspondence.} The exhaustion threshold \ensuremath{\nu *(R)} of
\S\,6.3 and its full symbolic expansion in Appendix A are derivations of
the structural form, not of the numerical value. The functional
dependences --- formation rate set by sampling rate, dissolution rate
set by per-level mortality plus predation plus starvation contributions,
exhaustion onset when the former drops below the latter --- are forced
by inherited machinery. The geometric proportionality constants
\ensuremath{G_{\mathrm{pred}}^{(k)}} and \ensuremath{G_{\mathrm{starv}}^{(k)},} together with the C(4,2) = 6
combinatorial factor inherited from Paper 3, are stated as worst-case
bounds whose exact refinement is reserved for Paper 5. The threshold's
qualitative behavior is forced by the derivation; its quantitative
pinning is not. The hierarchical lightspeed dispersion spectrum (\S\,7.3)
is structurally constrained by the upper bound but unbounded below; the
explicit \ensuremath{c^{(k)}} sequence depends on combinatorial enumeration deferred
to Paper 5.

\textbf{Predictions / Conjectures.} Falsifiable claims that follow from
the framework but whose closed forms are not derived within this paper.
The fully refined \ensuremath{\beta} = f(d, \ensuremath{K_{\psi,\mathrm{local}},} \ensuremath{K_{\psi,\mathrm{surround}},} \ensuremath{K_{\psi,\mathrm{subsystem}},}
k) of \S\,8.1 specifies the functional dependences across all five
arguments; f's explicit form is deferred. The cross-level relativistic
phenomena flagged in \S\,8.2 (variable apparent propagation, frame-relative
timing, embedded-medium dilation) are structural consequences of the
\ensuremath{c^{(k)}} dispersion and the field reinterpretation protocol; their
derivation as the substrate-level origins of relativistic phenomena
requires the closed-form \ensuremath{c^{(k)}} values and the reinterpretation
protocol of Paper 5. The identification of our measured speed of light
as \ensuremath{c^{(k)}} for some hierarchical level k determined by which level our
universe occupies (\S\,7.6) is a structural prediction whose empirical
content depends on Paper 5's simulation work.

\textbf{Deferred derivations carried into Paper 5.} Several results are
stated structurally in this paper and require Paper 5's simulation work
for closure: the closed-form \ensuremath{c^{(k)}} values at successive levels, the
exact integer \ensuremath{L_{p}^{(k)}} values, the numerical evaluation of \ensuremath{G_{\mathrm{pred}}^{(k)}}
and \ensuremath{G_{\mathrm{starv}}^{(k)},} the full closed form of f, the cross-level
relativistic derivation, the reinterpretation protocol that takes the
substrate from exhaustion to fundamental fields, the identification of
those fields and their initial conditions, and the Big Bang
correspondence. Paper 4 establishes the structural framework within
which these closures occur; Paper 5 supplies the closures. The
four-paper sequence of Volume 1 is the structural derivation; Paper 5 is
the numerical and protocol pinning that converts the structural results
into specific empirical content.

This four-way labeling applies throughout the paper. Every claim in the
body text falls into exactly one category.

\section*{Preface to Paper 4}
\addcontentsline{toc}{section}{Preface to Paper 4}

This paper is the fourth installment in Volume 1 of the Triadic
Measurement Substrate program. Volume 1, Paper 1 established the
geometric and ontological architecture of the substrate. Volume 1, Paper
2 developed the substrate's computational dynamics. Volume 1, Paper 3
established how computation emerges from informational substrate via
surround-driven repair, polarization-coherence bistability, recursive
coarse-graining, and the first computing subsystem. Paper 4 takes all
three as given and develops the macroscopic consequences promised at
Paper 3 \S\,6.5: the recursive hierarchies of subsystems, the conditions
under which a subsystem develops an internal model of its environment,
the dynamics of multiple subsystems sharing a meta-interstice
environment, the eventual exhaustion of accessible novelty as a
substrate region completes its configuration space, and the hierarchical
lightspeeds \ensuremath{c^{(k)}} developed in detail. As in Papers 1 through 3, no
new axioms are introduced. No new free numerical parameters are
introduced. The Principle of Generative Economy continues to govern
every selection.

Paper 4 is the connective tissue of Volume 1. Paper 3 derived the first
computing subsystem; Paper 5 will develop the reinterpretation protocol
by which a substrate region whose accessible configuration space has
exhausted is read under a new protocol, producing the fundamental fields
and initial conditions of a simulated universe. Between these two
endpoints lies the work of Paper 4: showing how subsystems multiply,
hierarchize, self-model, stress one another, diversify, and ultimately
exhaust their local configuration spaces --- the structural sequence
that makes reinterpretation both necessary and possible.

\begin{paperabstract}
Building on the first computing subsystem of Paper 3, we develop the
macroscopic consequences of subsystem dynamics. Six principal results
are derived. First, the recursive coarse-graining operator G of Paper 3
generalizes to \ensuremath{G^{(k)}} at hierarchical level k, with the
scale-separation requirements and the maximum hierarchical depth \ensuremath{k_{\mathrm{max}}}
for any bounded region following from the meta-level coherence
threshold. The dual recursion --- upward coarse-graining and downward
sub-subsystem nesting --- bounds the hierarchy above at \ensuremath{k_{\mathrm{max}}} and below
at the substrate-level minimum program extent \ensuremath{L_{p}.~\text{Second},} subsystems
grow via three structurally distinct modes --- annexation (extensive
growth via substrate absorption), imprinting (intensive growth via
correlation capture), and incorporation (intensive growth via structural
replication) --- each with its own derivation, mechanism, and ceiling.
Third, self-modeling is a structural property of every subsystem at
hierarchical level k \ensuremath{\geq} 1, not a special condition some subsystems happen
to meet. Every subsystem is therefore a measurement node, with the
triadic confirmation structure recurring at every level. Fourth, three
substrate-level stressors --- mortality (coherence failure), predation
(boundary annexation), and starvation (surround impoverishment) --- are
derived rigorously, each with a substrate-level mechanism and a
corresponding level-(k+1) signature at the meta-meta-confirmation that
reads the affected subsystem as a meta-tetrahedral corner. Fifth,
subsystem formation is overdetermined by three converging arguments:
Borel-Cantelli inevitability in the infinite-flop limit, per-event
Principle of Generative Economy preference, and flux-balance forcing
from signal conservation. Sixth, novelty exhaustion is derived as the
monotonic decay of the novelty-production rate \ensuremath{\nu (R,} \ensuremath{\tau )} once a
substantial fraction of the finite set \ensuremath{\Sigma (R)} of accessible stable
configurations has been sampled, with the exhaustion threshold \ensuremath{\nu *(R)}
given in closed symbolic form. The hierarchical lightspeed \ensuremath{c^{(k)}}
develops a dispersion spectrum across levels, bounded above by
\ensuremath{c^{(k+1)}} \ensuremath{\leq} \ensuremath{(1/\sqrt{2})} \ensuremath{c^{(k)}} and unbounded below, with each level's
measurement nodes observing the propagation speed of their own level. As
in Papers 1 through 3, every quantity is derived from the five axioms
and the Principle of Generative Economy; no new free numerical
parameters are introduced.

The arc of Paper 4 closes by naming the structural problem that Paper 5
inherits. At the exhaustion threshold of any bounded region, the
substrate holds informationally rich confirmed bits with no remaining
direct-read protocol capable of extracting novelty from them. Axiom 4
(causal anchoring) forbids overwriting the bits; Axiom 3 (signal
conservation) forbids losing the signals that produced them. The only
structurally available move is reinterpretation under a different
protocol --- the work of Paper 5.
\end{paperabstract}

\section*{I. Inheritance from Papers 1--3}
\addcontentsline{toc}{section}{I. Inheritance from Papers 1--3}

\subsection*{1.1 Inherited Structure}

Paper 4 takes the following structures from Papers 1 through 3 as given
without re-derivation. From Paper 1: the five axioms (Information
Fundamentalism, Measurement Imperative, Sphere as Primitive, Signal
Conservation, Causal Anchoring) and the Principle of Generative Economy
(PGE); the agent manifold \ensuremath{M\alpha} with agent spacing \ensuremath{\ell_{0}} \ensuremath{\equiv} \ensuremath{\ell_{p};} the FCC causal
stack with face-diagonal nearest-neighbor connectivity; the triadic
actualization rule \ensuremath{\bar{B}} = \ensuremath{3\bar{\alpha}_{2}} in 2D and \ensuremath{\bar{B}} = \ensuremath{4\bar{\alpha}_{3}} in 3D; the \ensuremath{1\to 4} multiplier
and informational closure; the binary alphabet \ensuremath{\bar{B}} \ensuremath{\in} \{0,1\} with
polarization \ensuremath{\chi} = \ensuremath{2\bar{B}} \ensuremath{-} 1; the flop as the atomic unit of state change
with no sub-flop time; the conjugate field structure \ensuremath{(\varphi ,} \ensuremath{\Pi )} and the wave
equation for \ensuremath{\varphi} at \ensuremath{c_{\mathrm{TMS}}} = \ensuremath{1/\sqrt{2};} the [[4,2,2]] structural
correspondence with logical qubit LQ1 (bit value) and LQ2
(causal-history inheritance); the polarization field \ensuremath{\psi} as parallel
propagating object on the same FCC lattice; the GUP prediction \ensuremath{\beta} = f(d,
\ensuremath{K_{\psi})} refined through Papers 2 and 3; and \ensuremath{U_{\mathrm{sh}}} as the maximum-entropy
substrate present at every depth of the causal stack.

From Paper 2: the polarization wave equation; the conjugate pair \ensuremath{(\psi ,}
\ensuremath{\Pi_{\psi});} the polarization-biased resolution rule \ensuremath{P(\bar{B}} = 1 \textbar{} \ensuremath{\Pi_{\psi})}
= (1 + \ensuremath{\Pi_{\psi})/2} and its Born-rule structural correspondence; the 3-input
majority gate; the universality result \{MAJ, constant\} \ensuremath{\to} \{AND, OR,
NOT\}; the flop operator T; and the formal program definition
\ensuremath{(T^{n}[C]} = C with K(C) \textless{} \textbar C\textbar).

From Paper 3: surround-driven repair and the Imprinting Theorem; the
imprint depth \ensuremath{d_{R}} = \ensuremath{\sqrt{2}} \ensuremath{\cdot} \ensuremath{L_{R}/\ell_{p};} the dissolution rate and effective
error-rate decay \ensuremath{(1/3)^{d};} the per-event error rate \ensuremath{p_{\mathrm{error}}} = (1 \ensuremath{-}
\textbar \ensuremath{\Pi_{\psi}}\textbar)/2; the coherence bistability threshold
\textbar \ensuremath{\Pi_{\psi}}\textbar\ensuremath{*(L_{R})} = 1 \ensuremath{-} \ensuremath{2\sqrt{c^{(k)}/(6}} \ensuremath{L_{R}));} the
coarse-graining operator G and the recursion to meta-level primitives;
the meta-level triadic rule, meta-level FCC structure, and meta-level
[[4,2,2]] stabilizer dynamics; the minimum program extent \ensuremath{L_{p};}
the \ensuremath{(C_{D},} \ensuremath{C_{O},} \ensuremath{C_{C})} trichotomy defining the subsystem; the minimum
subsystem extent \ensuremath{L_{\mathrm{sub}}} = \ensuremath{3L_{p}} + \ensuremath{\delta ;} the information-generation rate
differential between subsystem-hosting and static-program regions; the
hierarchical lightspeed formula \ensuremath{c^{(k)};} and the
selection-for-subsystems result.

These structures constitute the entirety of the foundation. Paper 4
introduces no further primitives.

\subsection*{1.2 The Questions Paper 4 Answers}

Paper 3 \S\,6.5 committed Paper 4 to five threads: hierarchies of
subsystems as nested computing configurations at multiple meta-scales;
conditions under which a subsystem develops an internal model of its
environment; dynamics of multiple subsystems sharing a meta-interstice
environment; eventual exhaustion of novelty as a substrate region
completes its accessible configuration space; and hierarchical
lightspeeds \ensuremath{c^{(k)}} developed in detail. These five threads are the
agreed scope and order of treatment, mapping to \S\,II--\S\,III, \S\,IV, \S\,V, \S\,VI,
and \S\,VII respectively.

Beyond the five committed threads, Paper 4 derives additional structural
results that emerge as the threads are developed: three growth modes
(annexation, imprinting, incorporation) with their respective ceilings;
three substrate-level stressors (mortality, predation, starvation) at
full proof with their two-level signatures; the inevitability of
subsystem formation via three converging arguments; inter-region and
intra-region diversity as derived consequences; and the region biography
\ensuremath{\mathcal{B}(R)} as a derived object whose phases organize the paper's spine. Each
emerges as a consequence of the inherited machinery, requiring no new
axioms or free parameters.

\section*{II. The Hierarchical Recursion}
\addcontentsline{toc}{section}{II. The Hierarchical Recursion}

\subsection*{2.1 The Recursive Coarse-Graining Operator \ensuremath{G^{(k)}}}

Paper 3 \S\,4.2 defined the coarse-graining operator G as the mapping from
a substrate-level polarization landscape on a region of extent \ensuremath{\geq} \ensuremath{L_{p}}
over a time window covering one program period to a single meta-level
state \ensuremath{\bar{P}_{p}} \ensuremath{\in} \{0, 1\} indexed by program identity p.~The construction
was scale-independent: the operator evaluated whether the configuration
satisfied T-closure plus stabilizer parity plus the coherence threshold,
and returned the meta-bit accordingly.

Paper 4 generalizes G to act at arbitrary hierarchical depth. The
operator \ensuremath{G^{(k)}} takes a configuration on a \ensuremath{\text{level}-(k-1)} region of
meta-extent \ensuremath{\geq} \ensuremath{L_{p}^{(k)}} over a time window covering one level-k flop
period \ensuremath{T_{p}^{(k)},} and returns a single level-k state \ensuremath{\bar{P}_{p}^{(k)}} \ensuremath{\in}
\{0, 1\} indexed by level-k program identity. The construction proceeds
by the same logic as Paper 3 \S\,4.2 applied at the new scale: identify
program-bearing regions at the relevant level, evaluate program
identity, output the meta-state. By the scale-independence of the axioms
(Paper 3 \S\,4.7), no new mechanism is required --- only the recursive
application of the same operator at the new scale.

The recursion is the geometric form of axiom scale-independence. The
five axioms and PGE do not refer to any particular scale; applied at
scale \ensuremath{\ell_{p}} they produce the substrate; applied at scale \ensuremath{L_{p}} they produce
the meta-substrate of programs; applied at scale \ensuremath{L_{p}^{(k)}} for any k
they produce the meta-substrate of level-k primitives. No new axiom, no
new free parameter, no new mechanism enters at any level.

The hierarchical-recursion structure derived here admits comparison with
an independent program developing a parallel architecture from different
starting premises. Rosas et al.~(2024) formalize hierarchical emergence
in complex systems via a lattice of nested computationally-closed levels
ordered by coarse-graining relationships. Their framework identifies
macroscopic levels as software-like processes that are causally and
informationally self-contained with respect to microscopic
instantiation, and develops the conditions under which such closure
obtains. The structural target --- a hierarchy of nested computational
levels generated by coarse-graining --- coincides with the recursion
produced here by \ensuremath{G^{(k)}.} The derivations differ in starting points:
the Rosas framework grounds hierarchical closure in
information-theoretic conditions on causal mechanisms, evaluating which
subsets of microscopic dynamics support self-contained macroscopic
computation. The TMS recursion grounds hierarchical closure in the
geometric scale-independence of the substrate's five axioms and PGE,
requiring no information-theoretic optimization at each level because
the same axioms force the same triadic confirmation structure wherever
they apply.

\subsection*{2.2 Scale-Separation Requirements}

For \ensuremath{G^{(k)}} to be well-defined, level-k structures must be cleanly
distinguishable from \ensuremath{\text{level}-(k-1)} noise. This requires two conditions:
\ensuremath{L_{p}^{(k)}} \ensuremath{\gg} \ensuremath{L_{p}^{(k-1)},} so that level-k primitives are spatially
distinguishable from \ensuremath{\text{level}-(k-1)} ones; and \ensuremath{T_{p}^{(k)}} \ensuremath{\gg} \ensuremath{T_{p}^{(k-1)},}
so that level-k flops are temporally distinguishable from \ensuremath{\text{level}-(k-1)}
ones.

These conditions are not postulated. They are forced by the coherence
threshold \textbar \ensuremath{\Pi_{\psi}}\textbar{}\emph{\ensuremath{(L_{R})} applied at the meta-level.
A candidate level-k program with spatial extent only marginally above
\ensuremath{L_{p}^{(k-1)}} would have insufficient room for its boundary conditions
to satisfy stabilizer parity against the local \ensuremath{\text{level}-(k-1)} noise; the
level-k stabilizer dynamics would read it as an unstable \ensuremath{\text{level}-(k-1)}
fluctuation rather than as a level-k primitive. The same balance between
per-flop noise injection and per-flop imprinting capacity that derived
\textbar \ensuremath{\Pi_{\psi}}\textbar{}}\ensuremath{(L_{R})} in Paper 3 \S\,3.3 forces clean
scale-separation at every level.

\subsection*{2.3 The Hierarchy Ceiling: \ensuremath{k_{\mathrm{max}}}}

For any bounded substrate region of spatial extent L and temporal extent
T (both measured in substrate units), the maximum hierarchical depth
\ensuremath{k_{\mathrm{max}}} is defined as the largest k satisfying \ensuremath{L_{p}^{(k)}} \textless{} L
and \ensuremath{T_{p}^{(k)}} \textless{} T. Beyond \ensuremath{k_{\mathrm{max}},} no level-k structure fits
within the region's spatial or temporal extent.

\ensuremath{k_{\mathrm{max}}} is finite for any finite (L, T). The hierarchy is geometrically
bounded above. The exact value of \ensuremath{k_{\mathrm{max}}} for a region depends on the
\ensuremath{L_{p}^{(k)}} and \ensuremath{T_{p}^{(k)}} ratios at each level, which depend on the
combinatorial enumeration of stabilizer-satisfying T-closed
configurations at that level --- a calculation deferred to Paper 5's
simulation work. Paper 4 establishes that \ensuremath{k_{\mathrm{max}}} exists and is finite;
the numerical value for a given region is a Paper 5 derivation.

\subsection*{2.4 The Russian Doll Floor: Downward Recursion}

Subsystems contain smaller subsystems at the same hierarchical level. A
subsystem of meta-extent \ensuremath{L_{\mathrm{sub}}^{(k)}} = \ensuremath{3L_{p}^{(k)}} + \ensuremath{\delta ^{(k)}} at
level k has interior substrate available for hosting \ensuremath{\text{level}-(k-1)}
sub-subsystems. The interior available extent is \ensuremath{L_{\mathrm{sub}}^{(k)}} minus the
boundary overhead occupied by the subsystem's three component programs.

The downward cascade --- sub-subsystems nesting inside
sub-sub-subsystems --- terminates when the interior of a level-k
subsystem can no longer accommodate even a single \ensuremath{\text{level}-(k-1)} subsystem.
Tracing this iteratively from the largest enclosing subsystem down: the
cascade terminates at level k = 0 (substrate programs) when \ensuremath{L_{p},} the
minimum program extent established in Paper 3 \S\,4.1, is reached. Below
\ensuremath{L_{p},} no programs exist; below \ensuremath{L_{p}^{(k)}} for any k, no level-k
subsystems exist. The Russian doll floor for any nesting chain is
therefore \ensuremath{L_{p}} at the substrate level.

A sub-sub-sub-subsystem chain can only descend until the innermost
subsystem has interior extent \textless{} \ensuremath{L_{p}.~\text{That}} is the geometric
floor on downward nesting, derived from the same minimum-program-extent
argument that bounds the substrate level.

\subsection*{2.5 The Dual Recursion and Its Asymmetry}

The two recursions --- upward (k = 0 \ensuremath{\to} k = 1 \ensuremath{\to} k = 2 \ensuremath{\to} \ldots) via
coarse-graining, and downward (k = N \ensuremath{\to} k = \ensuremath{N-1} \ensuremath{\to} \ldots{} \ensuremath{\to} k = 0) via
spatial nesting --- are distinct objects meeting at the same substrate.
Both have floors at \ensuremath{L_{p}} and ceilings at the region's spatial and
temporal extent. They are not, however, inverses of each other. The
asymmetry of the two operations is structural and matters for the
configuration count Paper 6 will use to bound regional novelty.

\subsection*{2.5.1 Lossy Coarse-Graining and Generative Nesting}

Coarse-graining is lossy. \ensuremath{G^{(k)}} takes a \ensuremath{\text{level}-(k-1)} configuration on
a region of meta-extent \ensuremath{\geq} \ensuremath{L_{p}^{(k)}} and returns a single bit
\ensuremath{\bar{P}^{(k)}_{p}} \ensuremath{\in} \{0, 1\}. The input is a complete polarization landscape,
potentially containing many thousands of \ensuremath{\text{level}-(k-1)} bits across the
time window. The output is one bit indexed by program identity. The
mapping is many-to-one: multiple distinct \ensuremath{\text{level}-(k-1)} landscapes map to
the same level-k \ensuremath{\bar{P}} value, because what gets preserved by the
coarse-graining is the program identity p, not the specific bit-by-bit
content of any implementation of that program.

This is forced by the program definition itself (Paper 2 \S\,6.2): a
program is a configuration causally closed under T with K(C) \textless{}
\textbar C\textbar. Multiple substrate-level configurations can
implement the same program identity --- they are members of the same
equivalence class under T-closure plus stabilizer parity plus
K-compressibility. \ensuremath{G^{(k)}} collapses the equivalence class to its
label. Two regions running identical majority-gate programs but with
different specific bit patterns coarse-grain to the same \ensuremath{\bar{P};} the
substrate-level distinction between the two implementations is lost at
the meta-level.

Downward nesting is generative. The recursion takes a level-k subsystem
and populates its interior with \ensuremath{\text{level}-(k-1)} sub-configurations. The
level-k subsystem itself does not specify the interior --- it specifies
only its boundaries (the three programs at meta-tetrahedral corners).
The interior substrate is available for sub-configurations but not
determined by them. When a \ensuremath{\text{level}-(k-1)} sub-subsystem forms inside the
level-k subsystem's interior, the \ensuremath{\text{level}-(k-1)} configuration is genuinely
new information that was not present at the level-k description.

So \ensuremath{G^{(k)}} and nesting are not inverses. The composition \ensuremath{G^{(k)}} \ensuremath{\circ}
nesting does not return the original configuration: the \ensuremath{\text{level}-(k-1)}
sub-configurations visible by nesting are absent from the level-k
coarse-grained description; they are generated by the nesting operation,
not stored. Similarly nesting \ensuremath{\circ} \ensuremath{G^{(k)}} does not recover what
coarse-graining lost: the specific implementation details of the
original \ensuremath{\text{level}-(k-1)} configuration that were absorbed into the
equivalence class are not retrievable from the level-k label alone.

The asymmetry is structural. Information lives at the level where it
lives. Going up the hierarchy loses \ensuremath{\text{level}-(k-1)} detail; going down the
hierarchy creates new \ensuremath{\text{level}-(k-1)} detail. Neither operation recovers
what the other does. This means the hierarchy is informationally rich
rather than trivial: each level has its own configuration space that is
not recoverable from the levels above or below. The region biography
\ensuremath{\mathcal{B}(R),} introduced below and developed across \S\,III--\S\,VII, must track each
level separately because each level is genuinely populated by
configurations the other levels do not see. The novelty argument of \S\,VI
depends on this per-level independence.

\subsection*{2.6 The Region Biography \ensuremath{\mathcal{B}(R)}}

A region biography is the time-indexed sequence of structural state of a
bounded region R as it traverses formation through saturation. Formally,

\emph{\ensuremath{\mathcal{B}(R)} = \ensuremath{\langle N_{\mathrm{sub}}^{(k)}(\tau ),} \{\ensuremath{K(C_{i})}\}\textbf{\{i \ensuremath{\in} R\}\ensuremath{(\tau ),}
\ensuremath{k_{\mathrm{max}}(R,} \ensuremath{\tau ),} \ensuremath{\Sigma_{\mathrm{sampled}}(R,} \ensuremath{\tau ),} \ensuremath{\nu (R,} \ensuremath{\tau )\rangle}}\{\ensuremath{\tau} \ensuremath{\in} \ensuremath{[\tau_{\mathrm{formation}},}
\ensuremath{\tau_{\mathrm{exhaustion}}]}\}}

where \ensuremath{N_{\mathrm{sub}}^{(k)}(\tau )} is the population of level-k subsystems in R at
flop time \ensuremath{\tau ;} \{\ensuremath{K(C_{i})}\}\ensuremath{_{i \in R}(\tau )} is the distribution of Kolmogorov
complexity values across configurations in R; \ensuremath{k_{\mathrm{max}}(R,} \ensuremath{\tau )} is the
current maximum hierarchical depth occupied in R (which rises over flop
time as subsystems form and coarse-grain into higher-level primitives);
\ensuremath{\Sigma_{\mathrm{sampled}}(R,} \ensuremath{\tau )} \ensuremath{\subseteq} \ensuremath{\Sigma (R)} is the set of stable configurations already
sampled by flop \ensuremath{\tau ;} and \ensuremath{\nu (R,} \ensuremath{\tau )} is the current rate at which novel (not
previously sampled) configurations enter \ensuremath{\Sigma_{\mathrm{sampled}}.}

Region biography is the object \S\,III--\S\,VII trace. \S\,III adds growth modes
(which change K and possibly the spatial extent of components). \S\,IV
characterizes self-modeling as a structural property of every subsystem
at k \ensuremath{\geq} 1. \S\,V adds multi-subsystem dynamics, stressors, and diversity
(which reshape the K-distribution and the subsystem population). \S\,VI
characterizes the saturation phase \ensuremath{(\Sigma_{\mathrm{sampled}}} \ensuremath{\to} \ensuremath{\Sigma (R)} and \ensuremath{\nu} \ensuremath{\to} 0). \S\,VII
develops the hierarchical lightspeed structure that conditions every
level's internal dynamics.

\ensuremath{\mathcal{B}(R)} is itself a configuration at the substrate level --- a time-indexed
record of structural state. By the recursion of \S\,2.1, \ensuremath{\mathcal{B}(R)} can be
coarse-grained to higher levels. A meta-subsystem can hold the biography
of its enclosed region as a configuration in its data register. This
biography-of-biographies setup is what Paper 5's reinterpretation
protocol will require: when the region's accessible configuration space
is exhausted via direct sampling, the biographical record itself becomes
the substrate for re-reading under a different protocol.

\section*{III. Growth Modes}
\addcontentsline{toc}{section}{III. Growth Modes}

\subsection*{3.1 Three Modes of Growth}

A subsystem can grow in three structurally distinct ways: by absorbing
adjacent substrate (annexation, extensive growth); by accumulating
compressible imprint structure within its existing extent (imprinting,
intensive correlation-capture growth); and by building internal copies
of external programs without absorbing or destroying them
(incorporation, intensive structural-replication growth). The three
modes are not alternatives. A single subsystem can grow via all three
simultaneously, on different axes, with different ceilings. Each mode is
derived from existing Paper 1--3 machinery: annexation from the
stabilizer parity condition applied at an expanding boundary, imprinting
from the Imprinting Theorem applied to a permeable data register,
incorporation from imprinting plus the additional condition of T-closure
on the recipient's substrate.

\subsection*{3.2 Annexation}

Annexation is the substrate-level mechanism by which a subsystem absorbs
adjacent substrate. Subsystem A's component (typically \ensuremath{C_{D}} or \ensuremath{C_{O})}
expands its spatial extent through a region R' that is either currently
unoccupied by any subsystem (raw substrate with \ensuremath{U_{\mathrm{sh}}-\text{like}} polarization)
or currently hosting a subsystem B whose coherence has dropped below its
threshold \textbar \ensuremath{\Pi_{\psi}}\textbar\ensuremath{*(L_{R}')} --- that is, dissolving anyway.

A's expansion preserves stabilizer parity at the new boundary because
A's interior imprints into the freed bits as the expansion proceeds. The
freed bits adopt A's polarization landscape via the same imprint
mechanism that originally formed A's boundary in Paper 3 \S\,II. B's
configuration, no longer maintained by its own components at the
dissolving boundary, falls below threshold and dissolves entirely. A
grows; B disappears.

Formally, A's component \ensuremath{C_{X}} (where X \ensuremath{\in} \{D, O, C\}) can expand into R'
if and only if the stabilizer parity condition
\textbar \ensuremath{\Pi_{\psi}}\textbar\ensuremath{_{(C_{X}}} \ensuremath{\cup} R') \ensuremath{\geq} \textbar \ensuremath{\Pi_{\psi}}\textbar\ensuremath{*(L(C_{X}} \ensuremath{\cup}
R')) is satisfied at every flop during the expansion. This is
geometrically restrictive: the larger \ensuremath{C_{X}} becomes, the higher the
coherence threshold required of its now-larger expanse, so unbounded
extensive growth fails when the threshold exceeds what A's interior can
sustain.

The annexation ceiling is therefore geometric. Extensive growth
saturates when either (a) the component's extent reaches the region's
extent, leaving no adjacent substrate to absorb, or (b) the coherence
threshold for the expanded component exceeds what the surround can
sustain. In typical regions of moderate \ensuremath{U_{\mathrm{sh}}} flux and ordinary subsystem
density, the second is the binding constraint: extensive growth halts
when the subsystem is too large to maintain coherence against its
surround.

In experiential framing reserved for Volume 2, annexation corresponds to
predation, feeding, and absorption. A subsystem eats by dissolving
another subsystem's structure and incorporating the freed bits into its
own components. The substrate-level mechanism is the directed expansion
of a component's boundary into adjacent substrate, preserving stabilizer
parity throughout.

\subsection*{3.3 Imprinting}

Imprinting is the substrate-level mechanism by which a subsystem's data
register accumulates compressible structure reflecting surround
regularities, without spatial extent changing. The Imprinting Theorem
(Paper 3 \S\,2.3) establishes that any boundary-permeable region imprints
its retarded surround history. When \ensuremath{C_{D}} is imprint-permeable --- that
is, its boundary with the surround communicates touch-vector signals
into its interior --- and its surround contains structured polarization
patterns, the imprint on \ensuremath{C_{D}'s} interior accumulates correlations
between \ensuremath{C_{D}'s} bits and surround features. Over many flops, \ensuremath{K(C_{D})}
rises because the imprinted pattern is compressible: the surround had
compressible structure to imprint, and that structure is captured rather
than created.

Formally, \ensuremath{K(C_{D})(\tau )} is monotonically non-decreasing in \ensuremath{\tau} for any
imprint-permeable \ensuremath{C_{D}} in a structured surround, with growth rate
proportional to the mutual information between \ensuremath{C_{D}'s} interior and the
surround's history within imprint depth \ensuremath{d_{D}.} The mutual information
measures how much of the surround's structure can be inferred from
\ensuremath{C_{D}'s} current interior --- and the imprinting mechanism maximizes this
iteratively, as established in Paper 3 \S\,II.

The imprint ceiling is informational. Intensive growth via imprinting
saturates when \ensuremath{K(C_{D})} approaches the algorithmic complexity upper bound
\textbar \ensuremath{C_{D}}\textbar{} \ensuremath{\cdot} log 2 established in Paper 3 \S\,3.1. Beyond this
bound, no additional compressible structure can fit in \ensuremath{C_{D}'s} bit budget
without spatial expansion. The data register is informationally full;
any further structure would require either growing \ensuremath{C_{D}'s} spatial extent
(annexation) or replacing existing imprint content with new content
(which is forbidden by causal anchoring as long as the existing content
remains coherent).

In experiential framing reserved for Volume 2, imprinting corresponds to
perception, observation, and sensory learning. The subsystem learns
about its surround by accumulating correlations --- without absorbing or
replicating anything. The model is statistical rather than structural:
\ensuremath{C_{D}'s} contents correlate with surround features, but the surround's
program-level structure is not internally reproduced.

\subsection*{3.4 Incorporation}

Incorporation is the substrate-level mechanism by which a subsystem's
data register builds an internal copy of an external program's structure
without absorbing the external program's substrate or destroying it.
Incorporation is imprinting plus an additional constraint.

\ensuremath{C_{D}} imprints from a surround that contains a program B --- a
configuration causally closed under T with \ensuremath{K(C_{B})} \textless{}
\textbar \ensuremath{C_{B}}\textbar. Ordinary imprinting (\S\,3.3) captures statistical
correlations between \ensuremath{C_{D}} and B's polarization landscape. Incorporation
occurs when the imprint reaches a state where \ensuremath{C_{D}'s} interior itself
satisfies T-closure on the imprinted pattern, with the imprinted region
within \ensuremath{C_{D}} being itself a program isomorphic in structure to B.

This is geometrically more restrictive than ordinary imprinting. The
imprint must capture not only B's polarization values but also satisfy
the T-closure conditions that make those values a program on \ensuremath{C_{D}'s}
substrate. This requires \ensuremath{C_{D}'s} interior to have enough room for B's
full structure \ensuremath{(L(C_{D})} \ensuremath{\geq} \ensuremath{L_{B})} and for the imprinted pattern to
satisfy the stabilizer parity conditions that make it a program rather
than a correlated pattern.

Formally, incorporation occurs at flop \ensuremath{\tau} when there exists a subregion
\ensuremath{C_{D}'} \ensuremath{\subseteq} \ensuremath{C_{D}} such that (a) \ensuremath{\psi} on \ensuremath{C_{D}'} matches \ensuremath{\psi} on B up to permutation
of bit positions, (b) \ensuremath{T^{n}[C_{D}']} = \ensuremath{C_{D}'} (T-closure on \ensuremath{C_{D}'s} own
substrate), and (c) \ensuremath{K(C_{D}')} = K(B). When all three hold, \ensuremath{C_{D}'} is a
program of identity \ensuremath{p_{B}} running on \ensuremath{C_{D}'s} substrate. B continues to
exist independently in the surround; \ensuremath{C_{D}} now contains a structural
replica of B operating on its own substrate.

The incorporation ceiling is structural. Each incorporated program of
identity \ensuremath{p_{B}} occupies space \ensuremath{L_{B}} in \ensuremath{C_{D}'s} interior, so the number of
distinct incorporated programs \ensuremath{C_{D}} can simultaneously hold is
approximately \ensuremath{L(C_{D})/L_{B}.} At higher hierarchical levels, this
constraint relaxes: a level-k \ensuremath{C_{D}} can incorporate \ensuremath{\text{level}-(k-1)} programs
more densely than substrate programs, because the \ensuremath{\text{level}-(k-1)} programs
have larger absolute spatial extent but occupy proportionally less of
\ensuremath{C_{D}'s} coarse-grained interior.

In experiential framing reserved for Volume 2, incorporation is the
substrate-level mechanism that corresponds to what is conventionally
called thinking with concepts. A subsystem with incorporated programs
holds working internal copies of external programs that can be applied
to other data. The conceptual framing is anthropomorphic and belongs to
Volume 2; the structural mechanism --- a region of \ensuremath{C_{D}} becomes a
program-hosting region in its own right, recursively --- is here.

\subsection*{3.5 The Sequencing of Growth Modes in Region Biography}

The three modes operate on different axes and require different inputs.
The dominance of each mode shifts across region biography because the
availability of each mode's required input changes as biography
progresses.

In the early phase of a region's biography, the region is freshly above
the static-program threshold. Many regions of raw substrate \ensuremath{(U_{\mathrm{sh}}-\text{like}}
polarization, no programs) are adjacent to newly-formed subsystems.
Annexation has abundant raw material. Imprinting and incorporation
require structured surrounds to draw from --- but the surrounds are
mostly unstructured at this stage. So annexation dominates by
availability of its raw material: subsystems primarily grow by consuming
the substrate around them.

In the middle phase, subsystems have multiplied. Raw substrate is
largely consumed; what remains is bordered by existing subsystems.
Adjacent regions are now populated by other subsystems whose
polarization landscapes are structured (each is a program with K(C)
\textless{} \textbar C\textbar). Annexation now requires breaking those
subsystems --- which is harder because they are coherent and
stabilizer-protected. The structured surrounds, however, are perfect for
imprinting: every neighbor is a compressible pattern that \ensuremath{C_{D}} can
capture as correlation. Imprinting dominates by availability of
structured surround.

In the late phase, imprint patterns have accumulated in many data
registers. Subsystems now contain compressed representations of their
neighbors. The next step --- building those representations as T-closed
programs internally --- is incorporation. The substrate for
incorporation is the imprints already present in \ensuremath{C_{D}'s} interior, which
only become available after sufficient imprinting has occurred.
Incorporation dominates last by availability of imprinted patterns to
elevate into internal programs.

The sequencing is forced by mode-specific input requirements: annexation
needs raw substrate, imprinting needs structured surround, incorporation
needs prior imprints. As biography progresses, the supply shifts from
raw to structured to imprinted, and the dominant mode shifts
accordingly. This is not strict sequencing --- all three modes can
operate at any phase --- but the dominant mode follows this trajectory
in regions whose biography starts at formation.

A region forming inside an already-populated surround skips ahead in the
sequence: if the surround is already structured when the region first
crosses the threshold, the early-phase annexation dominance is bypassed
and imprinting dominates from the start. The trajectory direction is the
same; the entry point varies. This nuance contributes to inter-region
diversity (\S\,V) and to the staggered nature of regional exhaustion (\S\,VI),
where regions of different biographical age are at different phases of
the sequence simultaneously.

\subsection*{3.6 Three Ceilings}

Each growth mode hits its ceiling for a different reason. Annexation has
a geometric ceiling: the coherence threshold rises with extent, and
growth halts when the threshold exceeds what the surround can sustain.
Imprinting has an informational ceiling: \ensuremath{K(C_{D})} cannot exceed
\textbar \ensuremath{C_{D}}\textbar{} \ensuremath{\cdot} log 2, beyond which no additional compressible
structure can fit in the existing bit budget. Incorporation has a
structural ceiling: the number of distinct incorporated programs is
bounded by \ensuremath{L(C_{D})/L_{B},} where \ensuremath{L_{B}} is the spatial extent of the
program being incorporated. The three different kinds of ceiling ---
geometric, informational, structural --- reflect that growth is
genuinely multi-axial: a subsystem can be limited in different ways on
different axes simultaneously, and the ceilings sometimes interact (a
subsystem at the informational ceiling can extend its budget by
extensive growth, raising the absolute K limit but inheriting the
geometric coherence-threshold constraint of its larger boundary).

\section*{IV. Self-Modeling as Structural Property}
\addcontentsline{toc}{section}{IV. Self-Modeling as Structural Property}

\subsection*{4.1 Self-Modeling Defined}

A subsystem self-models its environment when three conditions are
jointly satisfied. First, its data register \ensuremath{C_{D}} is imprint-permeable:
boundary touch-vector signals from the surround communicate into \ensuremath{C_{D}'s}
interior, and the Imprinting Theorem (Paper 3 \S\,2.3) governs the
resulting interior polarization landscape. Second, its operation kernel
\ensuremath{C_{O}} reads \ensuremath{C_{D}} as input: the substrate-level dynamics at the \ensuremath{(C_{D},}
\ensuremath{C_{O})} interface implement the 3-input majority gate computation (Paper 2
\S\,V) over \ensuremath{C_{D}'s} polarization values. Third, its control component \ensuremath{C_{C}}
selects operations in \ensuremath{C_{O}} based on the read result: the substrate-level
dynamics at the \ensuremath{(C_{O},} \ensuremath{C_{C})} and \ensuremath{(C_{D},} \ensuremath{C_{C})} interfaces propagate
\ensuremath{C_{O}'s} outputs into \ensuremath{C_{C}'s} operation-selection mechanism, closing the
loop from data through computation back to behavioral selection.

These three conditions are jointly what self-modeling requires.
Condition (a) gives correlation between \ensuremath{C_{D}} and surround. Condition (b)
makes \ensuremath{C_{D}'s} contents available to computation, distinguishing modeling
from passive imprinting. Condition (c) closes the loop, distinguishing
modeling from one-shot computation: the read result affects the
subsystem's subsequent operations, so the model is used.

\subsection*{4.2 Self-Modeling Is Structural}

Every subsystem at hierarchical level k \ensuremath{\geq} 1 satisfies the three
conditions of \S\,4.1 by its geometric construction. Self-modeling is not a
special condition that some subsystems meet and others miss --- it is a
structural property of subsystem-hood at any non-trivial hierarchical
depth.

By Paper 3 \S\,V, a subsystem at level k \ensuremath{\geq} 1 is the trichotomy \ensuremath{(C_{D},} \ensuremath{C_{O},}
\ensuremath{C_{C})} at meta-tetrahedral extent \ensuremath{L_{\mathrm{sub}}^{(k)}} = \ensuremath{3L_{p}^{(k)}} +
\ensuremath{\delta ^{(k)}.} Each component is itself a program at level k. The trichotomy
is geometrically arranged at three of the four meta-tetrahedral corners,
with the components' boundaries touching at the central meta-interstice.

Condition (a) is satisfied because \ensuremath{C_{D}'s} boundary with the surround is
imprint-permeable by the Imprinting Theorem: every boundary site of any
region of the FCC lattice is imprint-permeable, with imprint depth \ensuremath{d_{D}}
= \ensuremath{\sqrt{2}} \ensuremath{\cdot} \ensuremath{L_{D}/\ell_{p}.} There is no such thing as an imprint-impermeable
boundary in this substrate --- imprint-permeability is a structural
property of FCC connectivity.

Condition (b) is satisfied because the \ensuremath{(C_{D},} \ensuremath{C_{O})} interface is
precisely the meta-tetrahedral edge where \ensuremath{C_{D}'s} polarization signals
enter \ensuremath{C_{O}'s} parent set at the next meta-flop. The 3-input majority gate
of Paper 2 \S\,V.1 fires at this interface by the polarization-biased
resolution rule. The substrate-level dynamics at the interface implement
the read operation; no separate read mechanism is required.

Condition (c) is satisfied because \ensuremath{C_{C}} is constructed as the program
whose flop dynamics select among operation kernels (Paper 3 \S\,V.4). The
selection input to \ensuremath{C_{C}} is \ensuremath{C_{O}'s} output, propagated through the \ensuremath{(C_{O},}
\ensuremath{C_{C})} interface by the same touch-vector network that carries every
polarization signal in the substrate.

Therefore every minimum subsystem at k \ensuremath{\geq} 1 is a self-modeler. No
additional condition is required. Self-modeling is structural, not
conditional.

\subsection*{4.3 Subsystems Are Measurement Nodes}

A measurement node is a structure that reads its surround. By \S\,4.2,
every subsystem at k \ensuremath{\geq} 1 reads its surround through its data register.
Therefore every subsystem at k \ensuremath{\geq} 1 is a measurement node.

This identification holds at every hierarchical level by the recursion.
At k = 0, the agents of Paper 1 are measurement nodes at the substrate
level: three agents enclose one bit through triadic confirmation,
reading their immediate neighbors via the six touch vectors. At k = 1,
subsystems are measurement nodes at the meta-level: three programs
enclose one meta-bit at meta-tetrahedral corners, reading their
meta-surround through the data register. At k = 2, meta-subsystems are
measurement nodes at the meta-meta-level: three subsystems enclose one
meta-meta-bit, reading their meta-meta-surround. At every level, the
same triadic structure produces the measurement node at that level.

\ensuremath{\bar{B}} = \ensuremath{3\bar{\alpha}} is the bass note of every measurement event in TMS. Three
primitives enclose one bit at every scale. The substrate-level
confirmation, the meta-level confirmation, and every higher-level
confirmation are the same structural operation applied at different
scales. This recursive identification --- subsystem at level k equals
measurement node at level k --- is the substrate-level form of the
observer concept that active programs in quantum cosmology have been
developing in adjacent territory (Gorobey et al.~2025; Zhao, Harlow,
Usatyuk 2025). Those programs treat observers as bounded von Neumann
algebras within closed universes; TMS treats them as self-modeling
subsystems in bounded regions of the substrate. Both frameworks converge
on the same intuition that the observer is not external to the system
but a particular kind of internal structure with measurement properties
forced by its geometry.

\subsection*{4.4 Imprint Richness and Neighbor Density}

Self-modeling is structural, but the richness of the model varies with
surround content. A subsystem alone in its meta-interstice imprints the
three bounding programs that enclose it --- three static polarization
patterns at the meta-tetrahedral corners. The imprint is real but
sparse; \ensuremath{K(C_{D})} at saturation is low because the surround offered little
compressible structure beyond the three corner patterns.

A subsystem with neighboring subsystems in the same meta-interstice
imprints the bounding programs and the time-varying outputs of the
neighbors. Each neighbor contributes a dynamic polarization stream to
the surround. \ensuremath{C_{D}'s} imprint integrates these contributions, producing a
richer interior landscape --- higher \ensuremath{K(C_{D}),} more correlations
captured, more compressible structure made visible to \ensuremath{C_{O}} at the read
interface.

Formally, \ensuremath{K(C_{D})} at saturation is bounded above by \ensuremath{K(\psi_{\mathrm{surround}}} within
imprint depth \ensuremath{d_{D}).} A static surround caps \ensuremath{K(C_{D})} low because the
surround itself has low K; a dynamic, neighbor-rich surround permits
\ensuremath{K(C_{D})} to approach its informational ceiling \textbar \ensuremath{C_{D}}\textbar{} \ensuremath{\cdot}
log 2 because the surround contains enough varied compressible structure
to fill \ensuremath{C_{D}'s} bit budget. The same intensive growth mode (\S\,3.3)
produces qualitatively different outcomes depending on neighbor density.

\subsection*{4.5 Self-Modeling and Incorporation}

A subsystem whose imprint reaches incorporation-level structure (\S\,3.4)
holds internal programs that replicate external programs --- not just
statistical correlations. Self-modeling at this level is
structural-replicative: the subsystem does not merely have correlations
between \ensuremath{C_{D}} and its surround; it contains working internal copies of
programs it has observed in its surround. \ensuremath{C_{O}} can apply these internal
copies to data --- including other internal copies --- recursively.

This is the substrate-level form of conceptual representation. The
subsystem has not just an imprint of its environment but functional
internal programs that can be run, applied, and combined. In
experiential framing reserved for Volume 2, this is the substrate of
thinking with concepts. In Paper 4's structural framing, it is the upper
end of self-modeling capacity reachable by a subsystem at any given
hierarchical level.

\subsection*{4.6 Higher-Order Self-Modeling}

At k = 1, a subsystem self-models its meta-interstice surround --- the
bounding programs and any other subsystems sharing the same enclosure.

At k = 2, a meta-subsystem self-models its meta-surround, which contains
other meta-subsystems. The meta-surround can include the very subsystem
doing the modeling, if that subsystem is one component of its own
meta-environment. Self-modeling at k = 2 therefore includes the
possibility of modeling oneself as a component of one's environment.

This is forced by the recursion. A level-2 subsystem's data register
imprints from its level-2 surround, which is a configuration of level-1
subsystems. If the level-2 subsystem is itself composed of three level-1
subsystems (its components \ensuremath{C_{D}^{(2)},} \ensuremath{C_{O}^{(2)},} \ensuremath{C_{C}^{(2)}),} and
any of those level-1 subsystems is also a neighbor in the level-2
surround, then the level-2 imprint captures the modeling subsystem's own
components within its data register. The subsystem is modeling itself.

Higher k extends this. At level k, a subsystem can model itself as part
of an environment containing many \ensuremath{\text{level}-(k-1)} subsystems including its
own components. The depth of self-reference scales with hierarchical
depth: at higher k, the self-model includes more layers of nested
structure within the modeling subsystem itself.

Self-modeling-of-self emerges at k \ensuremath{\geq} 2 as a derived structural
consequence of the recursion, with no new mechanism. The experiential
reading of self-modeling-of-self --- what it is to be such a subsystem
from inside --- is reserved for Volume 2, which treats the agent-side
complement of the measurement-node structure derived here.

\subsection*{4.7 Synthesis}

Four claims tie together. Self-modeling is structural at every k \ensuremath{\geq} 1:
any minimum subsystem self-models by its geometric construction (\S\,4.2).
Every such self-modeling subsystem is a measurement node at its level:
subsystem and measurement node are the same object structurally (\S\,4.3).
The richness of a subsystem's model scales with the K-complexity of its
surround, set by neighbor density: same structure, different content
(\S\,4.4). At the upper end, self-modeling reaches structural-replicative
form when imprints become T-closed internal programs (\S\,4.5, connected to
\S\,3.4). At hierarchical depth k \ensuremath{\geq} 2, the self-modeling capacity includes
self-modeling-of-self, the substrate of higher-order self-reference
(\S\,4.6). The whole architecture follows from the inherited \ensuremath{(C_{D},} \ensuremath{C_{O},}
\ensuremath{C_{C})} trichotomy plus the Imprinting Theorem, without any new axiom or
mechanism.

\section*{V. Multi-Subsystem Dynamics, Stressors, and Diversity}
\addcontentsline{toc}{section}{V. Multi-Subsystem Dynamics, Stressors, and Diversity}

\subsection*{5.1 Inevitability of Subsystem Formation}

Before deriving multi-subsystem dynamics, we establish that subsystems
form. Subsystem formation in a region above the static-program threshold
is overdetermined by three converging arguments: inevitable in the
infinite-flop limit, preferred at every selection event, and forced by
signal conservation.

The first argument is Borel-Cantelli. \ensuremath{U_{\mathrm{sh}}} continuously samples all
configurations in any region (Paper 3 \S\,3.5); subsystem configurations
are a non-empty subset of stable T-closed configurations (the worked
example of Paper 3 \S\,5.4 provides an explicit minimum subsystem); the
per-event probability of sampling a subsystem-bearing configuration is
non-zero. Over infinite flop time, any event with non-zero per-event
probability occurs almost surely. Therefore subsystem configurations are
sampled with probability approaching 1 over finite expected flop time in
any region with above-threshold coherence.

The second argument is PGE preference. The Principle of Generative
Economy selects minimum sufficient configurations that maximize
generative potential. In a region above the static-program threshold,
subsystem configurations have higher generative capacity than static
programs because they transform patterns rather than merely persisting
them, by the universality result of Paper 2 \S\,V.3. At every selection
event --- not just statistically over infinite time --- PGE prefers
subsystems. The first argument gives inevitability in the limit; this
argument gives per-event preference.

The third argument is flux-balance forcing. \ensuremath{U_{\mathrm{sh}}} delivers signal flux to
a region at a rate proportional to the region's volume \ensuremath{V_{R}.} Static
programs consume flux at a rate proportional to their boundary area \ensuremath{A_{R}}
times their imprint capacity, which scales with boundary area only. The
flux balance for a region containing only static programs is generally
not satisfied: \ensuremath{V_{R}-\text{scaled}} delivery exceeds \ensuremath{A_{R}-\text{scaled}} consumption.
Axiom 3 (signal conservation) requires some configuration to consume the
excess; signals cannot accumulate without being expressed in
confirmations. Subsystems consume flux at a higher per-volume rate than
static programs because each transformation event consumes more signals
than a persistence event. Flux balance therefore forces subsystem
density to rise until consumption equals delivery.

The three arguments converge: inevitable in the infinite-time limit
(Borel-Cantelli), preferred at every selection event (PGE), forced by
signal conservation (flux-balance). Subsystem formation is
overdetermined. In any region above the static-program threshold but
below the subsystem-saturation threshold (defined in \S\,VI), subsystems
must form.

\subsection*{5.2 Meta-Meta-Confirmation}

When three subsystems at level k occupy three corners of a level-(k+1)
meta-tetrahedron, their joint outputs confirm a level-(k+1) P-state
\ensuremath{\bar{P}^{(k+1)}} at the fourth corner. This is the exact recursion of
substrate-level confirmation: three measurement nodes at
meta-tetrahedral corners enclose one bit. The same triadic minimum,
applied at level k+1.

The polarization inputs at the meta-meta-confirmation are the outputs of
the three subsystems --- the result of each subsystem's per-meta-flop
computation, propagated to the meta-tetrahedral edge that subsystem
shares with the meta-interstice. The polarization-biased resolution rule
from Paper 2 \S\,IV applies unchanged: \ensuremath{P(\bar{P}^{(k+1)}} = 1) = (1 +
\ensuremath{\Pi_{\psi}^{(k+1)})/2,} where \ensuremath{\Pi_{\psi}^{(k+1)}} is the average polarization of the
three subsystem outputs at the meta-tetrahedral edges feeding the
meta-interstice.

By the universality result of Paper 2 \S\,V.3, the meta-meta-confirmation
implements the 3-input majority gate over subsystem outputs. The
computational universality of the substrate is preserved at every
hierarchical level. Three subsystems form one meta-meta-confirmation in
the same way three programs form one meta-confirmation in the same way
three agents form one substrate-bit confirmation: \ensuremath{\bar{B}} = \ensuremath{3\bar{\alpha}} as the
recurring bass note.

\subsection*{5.3 Coupling Timescale Constraint}

For three subsystems at level k to coherently confirm a meta-state at
level k+1, each subsystem's output must reach the meta-interstice at the
moment of meta-meta-confirmation. Outputs propagate at the
meta-propagation speed \ensuremath{c^{(k)}} = \ensuremath{(L_{p}(k)/L_{p}(k-1))} \ensuremath{\cdot}
\ensuremath{(T_{p}(k-1)/T_{p}(k))} \ensuremath{\cdot} c (Paper 3 \S\,6.1).

Imprint depth for the three subsystem outputs to reach the
meta-interstice corner from the meta-tetrahedral edges is \ensuremath{d_{\mathrm{imprint}}} =
\ensuremath{\sqrt{2}} \ensuremath{\cdot} \ensuremath{L_{p}(k+1)/L_{p}(k)} measured in level-k flops. For
meta-meta-confirmation to occur coherently, the level-(k+1) flop period
must satisfy \ensuremath{T_{p}^{(k+1)}} \ensuremath{\geq} \ensuremath{d_{\mathrm{imprint}}} \ensuremath{\cdot} \ensuremath{T_{p}^{(k)},} which yields

\emph{\ensuremath{T_{p}^{(k+1)}} \ensuremath{\geq} \ensuremath{\sqrt{2}} \ensuremath{\cdot} \ensuremath{(L_{p}}\textbf{\ensuremath{(k+1)/L_{p}}}(k)) \ensuremath{\cdot} \ensuremath{T_{p}^{(k)}}}

This is a derived constraint on the hierarchical flop periods.
Substituting into the \ensuremath{c^{(k+1)}} formula yields a derived upper bound
\ensuremath{c^{(k+1)}} \ensuremath{\leq} \ensuremath{(1/\sqrt{2})} \ensuremath{\cdot} \ensuremath{c^{(k)}} in the tight-coupling regime --- a result
developed further in \S\,VII when the hierarchical lightspeed structure is
treated in detail. The constraint here establishes that the upward
recursion is not free: the geometry of meta-tetrahedral confirmation
forces the flop periods to grow with the meta-spacing ratios at each
level.

\subsection*{5.4 Three Stressors: Setup}

A stressor is any substrate-level dynamic that drives a subsystem below
the coherence threshold \textbar \ensuremath{\Pi_{\psi}}\textbar\ensuremath{*(L_{R})} or otherwise
disrupts its persistence. Three principal stressors emerge from the
substrate's own machinery without any additional mechanism: mortality
(internal coherence failure), predation (boundary annexation by an
active neighbor), and starvation (impoverishment of the structured
surround that maintains the subsystem's coherence). Each is derived at
two levels: a substrate-level mechanism that drives the stressor, and a
level-(k+1) signature visible at the meta-meta-confirmation that reads
the affected subsystem as a meta-tetrahedral corner.

\subsection*{5.5 Mortality}

Substrate-level mechanism: By Paper 3 \S\,3.4, a region's coherence
\textbar \ensuremath{\Pi_{\psi}}\textbar\ensuremath{_{R}} must exceed the threshold
\textbar \ensuremath{\Pi_{\psi}}\textbar\ensuremath{*(L_{R})} for the region to sustain program-bearing
content. A subsystem maintains its coherence through three sources:
internal stabilizer repair (Paper 1 \S\,V), surround imprint of structured
neighbors (Paper 3 \S\,II), and the per-flop polarization output of its own
operation kernel. When any of these sources degrades --- internal
stabilizer parity failing due to accumulated uncorrectable errors,
surround structure dissolving and imprint becoming noisy, operation
kernel output dropping below coherent --- the subsystem's coherence
falls below threshold. Once below threshold, Paper 3 \S\,3.4 establishes
that the configuration dissolves over imprint depth \ensuremath{d_{R}} = \ensuremath{\sqrt{2}} \ensuremath{\cdot}
\ensuremath{L_{\mathrm{sub}}/\ell_{p}} flops into \ensuremath{U_{\mathrm{sh}}-\text{like}} noise.

Mortality is the substrate-level event of falling below threshold and
dissolving. No new mechanism is required; it is the negative of the
persistence condition Paper 3 already established. A subsystem dies when
its coherence sources collectively fall below the threshold required to
maintain its structural integrity.

Level-(k+1) signature: When subsystem B at one corner of a
meta-tetrahedral confirmation dissolves, its output polarization stream
at the corner goes silent --- the corner stops contributing a structured
signal to the meta-meta-confirmation. The meta-meta-confirmation now
reads only two coherent inputs plus one effectively zero input.
Formally, \ensuremath{\Pi_{\psi}^{(k+1)}} at the affected meta-interstice drops from \ensuremath{(\chi_{A}}
+ \ensuremath{\chi_{B}} + \ensuremath{\chi_{C})/3} to approximately \ensuremath{(\chi_{A}} + 0 + \ensuremath{\chi_{C})/3,} where the 0
represents B's silenced corner. The meta-meta-confirmation's resolution
probability shifts toward 1/2 --- the confirmation becomes effectively a
coin flip rather than a biased computation. At higher hierarchical
levels, this propagates as a degradation of the meta-subsystem's
computational coherence: an event at level k is the structural correlate
of what experiential framing reserved for Volume 2 reads as absence at
level k+1, the silence of the corner registered by the larger pattern.

\subsection*{5.6 Predation}

Substrate-level mechanism: Subsystem A's component (typically \ensuremath{C_{D}} or
\ensuremath{C_{O})} expands its extent through annexation (\S\,3.2) into a region
currently occupied by subsystem B. A's expansion preserves stabilizer
parity at the new boundary because A's interior imprints into the freed
bits as the expansion proceeds. B's configuration, no longer maintained
by its own components which are being dissolved at the boundary, falls
below threshold and dissolves. A grows; B disappears.

Predation is distinct from mortality. In mortality, B fails on its own
--- its coherence sources fail without an active aggressor. In
predation, A actively dissolves B by annexing B's substrate. The
substrate-level signature is the directed expansion of A's boundary into
B's region, with A's interior imprint replacing B's polarization
landscape.

Level-(k+1) signature: B's corner of the meta-tetrahedral confirmation
is now occupied by A's pattern, since A's component has expanded into
the substrate that B previously occupied. \ensuremath{\Pi_{\psi}^{(k+1)}} at the
meta-interstice is approximately \ensuremath{(\chi_{A}} + \ensuremath{\chi_{A}} + \ensuremath{\chi_{C})/3} = \ensuremath{(2\chi_{A}} +
\ensuremath{\chi_{C})/3.} The signature is polarization homogenization: two corners carry
the same pattern (A's), where they previously carried distinct patterns
(A's and B's). The meta-meta-confirmation reads reduced diversity at
this corner-set, and the information-carrying capacity of the
meta-tetrahedral confirmation degrades --- three distinct inputs
collapsed to two effectively-identical inputs plus one distinct one.

\subsection*{5.7 Starvation}

Substrate-level mechanism: A subsystem's coherence depends on imprint
from a structured surround (\S\,4.4). When the surround becomes
uninformative --- neighbors dissolve, or surrounding programs stabilize
into low-K static patterns --- the subsystem's \ensuremath{C_{D}} imprints noise
rather than structure. Formally, when \ensuremath{K(\psi_{\mathrm{surround}}} within imprint
depth) approaches log(volume), the imprinted \ensuremath{C_{D}} has \ensuremath{K(C_{D})}
approaching log(\textbar \ensuremath{C_{D}}\textbar), the trivial floor. \ensuremath{C_{D}'s}
content becomes algorithmically uninteresting; \ensuremath{C_{O}'s} operations on a
uniform \ensuremath{C_{D}} produce uniform outputs; the subsystem's operation kernel
cannot perform non-trivial computation because there is no structured
data to process. The subsystem persists structurally (its stabilizer
parity may remain intact internally) but its computational function goes
idle for lack of structured input.

Starvation is distinct from mortality and predation. In mortality, B
fails internally. In predation, A actively destroys B. In starvation, B
persists structurally but stops computing because its environment no
longer feeds it structured input. Predation requires predators with
active annexation dynamics; mortality results from internal failure;
starvation requires environmental impoverishment. The three stressors
arise from genuinely different substrate-level dynamics.

Over time, without dynamic output, the subsystem's contribution to
neighboring imprints also falls to noise. This creates a positive
feedback: starved subsystems contribute uninformative surrounds to their
neighbors, who in turn starve. Starvation propagates through neighbor
networks more slowly than predation but more cumulatively.

Level-(k+1) signature: B's output at the meta-tetrahedral corner becomes
constant or low-K. \ensuremath{\Pi_{\psi}^{(k+1)}} at the meta-interstice includes a
corner contributing a frozen value rather than a dynamic stream. The
meta-meta-confirmation continues, but its variability is reduced ---
fewer distinct meta-meta-outcomes are reached because one input source
is informationally stuck. The signature is output rigidification: a
corner contributing a fixed or near-fixed value over many
meta-meta-flops, where it previously contributed dynamic variation.

\subsection*{5.8 Stressor Interactions and Selection}

The three stressors interact. Mortality is a passive failure mode;
predation is an active mechanism that can cause mortality in the prey;
starvation is a slow degradation that can lead to mortality if it
persists. Predation requires predators with active annexation dynamics.
Starvation often results from a populated region exhausting its raw
substrate, leaving only structured neighbors whose outputs may become
rigidified at long times. Mortality is the terminal phase of either
active stressor, as well as a possible terminal phase of subsystems that
simply degrade from internal accumulated errors.

Subsystems exposed to these stressors face a selection environment that
is not externally imposed but generated by the substrate's own
multi-subsystem dynamics. A subsystem persists by avoiding all three:
maintaining internal coherence against mortality, maintaining defensive
boundary configurations against predation, and maintaining access to
structured surround against starvation. The configurations that persist
are those whose internal structure and behavior happen to navigate all
three simultaneously.

\subsection*{5.9 Diversity: Two Mechanisms}

Diversity of subsystem configurations across the substrate has two
independent origins.

Inter-region diversity: Different regions of the substrate have
different local conditions. Different \ensuremath{U_{\mathrm{sh}}} disruption rates, different
stressor profiles (some regions have abundant raw substrate and low
predation pressure; others have dense subsystem populations and high
predation pressure), different surround histories that imprint different
patterns. The subsystems that survive in different regions are biased by
different selection pressures. The population of surviving subsystems
across the substrate is therefore diverse across regions --- different
stressor profiles produce different surviving subsystem distributions.

Intra-region diversity: Within a single region under a single stressor
profile, the survival set \ensuremath{\Sigma_{\mathrm{survive}}(R,} S) --- the set of subsystem
configurations that satisfy the survival conditions under stressor
profile S --- is generically large. A single stressor (e.g., the need to
obtain structured input under a partial starvation pressure) admits many
distinct solutions: active acquisition of structured neighbors via
boundary imprint exchange, passive imprint of stable bounding programs,
internal generation of structure via incorporation of remembered
patterns, restructuring the operation kernel to extract more from less.
The continuous-sampling mechanism (Paper 3 \S\,3.5) populates \ensuremath{\Sigma_{\mathrm{survive}}(R,}
S) with whatever configurations \ensuremath{U_{\mathrm{sh}}} instantiates that satisfy the
survival conditions. Multiple distinct configurations coexist within the
same survival set.

Intra-region diversity is the formal answer to the question of why
diverse creatures emerge in the same environment under the same
stressors. The stressor underdetermines the solution; the substrate
samples the underdetermined solution space; multiple solutions persist.
This complements inter-region diversity: even before regions are
compared, each region's internal subsystem population is itself diverse
because the local stressor profile underdetermines the solutions to its
own pressures.

\section*{VI. Novelty Exhaustion}
\addcontentsline{toc}{section}{VI. Novelty Exhaustion}

\subsection*{6.1 Finiteness of Accessible Configuration Space}

Let R be a bounded substrate region of spatial extent L and temporal
extent T in substrate units. Let \ensuremath{\Sigma_{k}(R)} be the set of stable
configurations at hierarchical level k that R can host ---
configurations causally closed under \ensuremath{T^{(k)},} satisfying stabilizer
parity at level k, with K(C) \textless{} \textbar C\textbar. Then
\ensuremath{\Sigma_{k}(R)} is finite for every k \ensuremath{\leq} \ensuremath{k_{\mathrm{max}}(R),} and \ensuremath{\Sigma (R)} =
\ensuremath{\bigcup _{k=0}^{k_{\mathrm{max}}}} \ensuremath{\Sigma_{k}(R)} is finite.

\begin{proof}
At level k, the number of bit positions in R is bounded by
\textbar R\textbar\ensuremath{/L_{p}^{(k)3}.} The total number of binary
configurations of those bits is bounded by
2(\textbar R\textbar\ensuremath{/L_{p}(k)^{3}).} The stable subset \ensuremath{\Sigma_{k}(R)} is a strict
subset of this total: most binary configurations do not satisfy
\ensuremath{T^{(k)}-\text{closure}} plus stabilizer parity plus K \textless{}
\textbar C\textbar. So \ensuremath{\Sigma_{k}(R)} is finite at every k. Summing over k from
0 to \ensuremath{k_{\mathrm{max}}(R),} \ensuremath{\Sigma (R)} is finite as a finite sum of finite sets.

The per-level structure matters. By \S\,2.5.1 (the asymmetry of
coarse-graining and nesting), configurations at different levels are
informationally independent: coarse-graining loses \ensuremath{\text{level}-(k-1)} detail,
nesting generates new \ensuremath{\text{level}-(k-1)} detail. So \ensuremath{\Sigma_{k}(R)} at each level
contributes independent novelty. \ensuremath{\Sigma (R)} is genuinely the sum, not a
relabeling of a single level's contribution. A region with \ensuremath{k_{\mathrm{max}}} = 1
has fewer accessible configurations than a region with \ensuremath{k_{\mathrm{max}}} = 5,
because the latter has five independent configuration spaces to sample
(substrate plus four higher levels) rather than one.

The per-level independence delays exhaustion. The hierarchy is not just
structural decoration on top of a single configuration count; it
genuinely multiplies the accessible novelty available to the region over
its biographical history.

\end{proof}

\subsection*{6.2 Novelty Production Rate}

Define \ensuremath{\nu (R,} \ensuremath{\tau )} as the rate of new (not previously sampled) stable
configurations entering \ensuremath{\Sigma_{\mathrm{sampled}}(R,} \ensuremath{\tau )} per unit substrate flop time at
flop \ensuremath{\tau .} Three factors determine \ensuremath{\nu :} the size of \ensuremath{\Sigma (R),} the fraction of
\ensuremath{\Sigma (R)} already sampled, and the instantaneous sampling rate set by \ensuremath{U_{\mathrm{sh}}}
flux and the continuous-sampling mechanism of Paper 3 \S\,3.5.

Formally, the expected \ensuremath{\nu (R,} \ensuremath{\tau )} at flop \ensuremath{\tau} is approximately

\emph{\ensuremath{\nu (R,} \ensuremath{\tau )} \ensuremath{\approx} \ensuremath{\rho_{\mathrm{sample}}(R)} \ensuremath{\cdot} (\textbar \ensuremath{\Sigma (R)}\textbar{} \ensuremath{-}
\textbar \ensuremath{\Sigma_{\mathrm{sampled}}(R,} \ensuremath{\tau )}\textbar) / \textbar \ensuremath{\Sigma (R)}\textbar{}}

where \ensuremath{\rho_{\mathrm{sample}}(R)} is the substrate's instantaneous sampling rate in R,
roughly constant for a given region under steady \ensuremath{U_{\mathrm{sh}}} flux.

When \textbar \ensuremath{\Sigma_{\mathrm{sampled}}(R,} \ensuremath{\tau )}\textbar{} is small relative to
\textbar \ensuremath{\Sigma (R)}\textbar, the second factor is near 1 and \ensuremath{\nu} is high. As
\textbar \ensuremath{\Sigma_{\mathrm{sampled}}(R,} \ensuremath{\tau )}\textbar{} approaches \textbar \ensuremath{\Sigma (R)}\textbar,
the second factor approaches 0 and \ensuremath{\nu} decays monotonically. This is not a
postulate; it is a counting argument. Every novel configuration is
sampled at most once before becoming non-novel. The pool of un-sampled
configurations shrinks monotonically as biographical history
accumulates, and the rate of novelty production decays with the pool.

\subsection*{6.3 The Exhaustion Threshold}

Define \ensuremath{\nu *(R)} as the critical novelty rate below which subsystem
formation can no longer keep pace with stressor-driven dissolution in R.

At \ensuremath{\nu} \textgreater{} \ensuremath{\nu}\emph{, new subsystems form (via inevitability,
\S\,5.1) faster than existing subsystems dissolve (via stressors,
\S\,5.4--\S\,5.7). Subsystem density rises or stabilizes. At \ensuremath{\nu}
\textbf{\textless{}} \ensuremath{\nu}}, new subsystem formation slows because the
sampling pool is depleted --- \ensuremath{U_{\mathrm{sh}}} keeps sampling, but most samples now
hit configurations already in \ensuremath{\Sigma_{\mathrm{sampled}},} which provide no new
structure. Meanwhile, stressors continue to dissolve existing subsystems
at undiminished rates. Net subsystem density falls.

\ensuremath{\nu *(R)} is set by the region's stressor profile. The closed-form symbolic
expression follows from balancing formation rate against dissolution
rate:

\emph{\ensuremath{\nu}}(R) = \textbar \ensuremath{\Sigma (R)}\textbar{} \ensuremath{\cdot} \ensuremath{\Sigma_{k}} \ensuremath{[d_{\mathrm{mort}}^{(k)}} +
\ensuremath{d_{\mathrm{pred}}^{(k)}} + \ensuremath{d_{\mathrm{starv}}^{(k)}]} / \ensuremath{(\rho_{\mathrm{sample}}(R)} \ensuremath{\cdot} \ensuremath{P_{\mathrm{coherent}})*}

where \ensuremath{d_{\mathrm{mort}}^{(k)},} \ensuremath{d_{\mathrm{pred}}^{(k)},} \ensuremath{d_{\mathrm{starv}}^{(k)}} are the level-k
dissolution rates derived in \S\,5.4--\S\,5.7, \ensuremath{\rho_{\mathrm{sample}}(R)} is the substrate
sampling rate, and \ensuremath{P_{\mathrm{coherent}}} is the probability that a sampled
configuration meets the coherence threshold
\textbar \ensuremath{\Pi_{\psi}}\textbar\ensuremath{*(L_{\mathrm{sub}}^{(k)}).} Each dissolution rate has a
derived form (worked out in Appendix A) that depends on the per-event
error rate \ensuremath{p_{\mathrm{error}}} = (1 \ensuremath{-} \textbar \ensuremath{\Pi_{\psi}}\textbar)/2, the
[[4,2,2]] code distance, and the local subsystem density. The
combinatorial factors in the dissolution-rate expressions (notably the
C(4,2) = 6 upper bound on per-tetrahedral-unit two-error probability)
are flagged as worst-case upper bounds whose exact numerical refinement
is reserved for Paper 5's simulation work.

The exhaustion threshold \ensuremath{\tau}\emph{(R) is the flop time at which \ensuremath{\nu (R,} \ensuremath{\tau )}
first drops below \ensuremath{\nu}}(R). After \ensuremath{\tau}\emph{, the region cannot sustain its
previous subsystem density. Subsystem populations decline; hierarchy
depths collapse downward as upper levels lose component subsystems; the
region's biography enters its terminal phase. Aggressive predation
environments have high \ensuremath{\nu}} (subsystems are destroyed rapidly, requiring
fast novelty influx); quiescent environments have low \ensuremath{\nu *.}

\subsection*{6.4 Region Biography Fully Formalized}

The region biography \ensuremath{\mathcal{B}(R)} introduced in \S\,2.6 is now fully specified by
combining the structural development of \S\,III--\S\,V with the novelty
dynamics of \S\,6.1--\S\,6.3. The biography traverses three phases.

Phase I --- Formation and Growth \ensuremath{(\tau} \ensuremath{\in} \ensuremath{[\tau_{\mathrm{formation}},} \ensuremath{\tau_{\mathrm{mid}}]):}
Subsystems form by the inevitability argument of \S\,5.1. They grow via the
three modes of \S\,III, with annexation dominating early when raw substrate
is abundant, transitioning to imprinting dominance as subsystems
multiply and surrounds become structured. Sub-subsystems nest downward
(\S\,2.4); coarse-grained meta-subsystems form upward (\S\,2.1). \ensuremath{N_{\mathrm{sub}}^{(k)}}
rises at every level; the K-distribution shifts toward higher values as
data registers accumulate compressible imprints; \ensuremath{k_{\mathrm{max}}(R,} \ensuremath{\tau )} climbs as
more hierarchical levels are populated.

Phase II --- Maturity \ensuremath{(\tau} \ensuremath{\in} \ensuremath{[\tau_{\mathrm{mid}},} \ensuremath{\tau *(R)]):} \ensuremath{\Sigma_{\mathrm{sampled}}} approaches
saturation but \ensuremath{\nu (R,} \ensuremath{\tau )} \textgreater{} \ensuremath{\nu *(R),} so the region still
produces novel configurations faster than stressors dissolve existing
subsystems. Subsystem density is stable; hierarchy is at or near \ensuremath{k_{\mathrm{max}};}
incorporation-mode growth dominates (\S\,3.5) as imprinted patterns are
elevated into internal programs. Stressor dynamics generate intense
diversity selection (\S\,5.8--\S\,5.9), producing both inter-region and
intra-region diversity. This is the phase of richest configuration
density per unit substrate.

Phase III --- Exhaustion \ensuremath{(\tau} \textgreater{} \ensuremath{\tau}\emph{(R)): \ensuremath{\nu (R,} \ensuremath{\tau )}
\textbf{\textless{}} \ensuremath{\nu}}(R). New formation lags dissolution.
\ensuremath{N_{\mathrm{sub}}^{(k)}} declines; \ensuremath{k_{\mathrm{max}}} drops as upper-level meta-subsystems lose
component subsystems and decoalesce; K-distribution drifts toward
already-sampled patterns; the region is no longer producing
computationally interesting novelty. Configurations that exist continue
to operate, but no new compressible structure enters the region's
biographical record.

\ensuremath{\mathcal{B}(R)} is itself a configuration at the substrate level --- a time-indexed
record of structural state. By the recursion of \S\,2.1, \ensuremath{\mathcal{B}(R)} can be
coarse-grained to higher levels. A meta-subsystem can hold the biography
of its enclosed region as a configuration in its data register. This
biography-of-biographies setup is the substrate-level prerequisite for
the work of Paper 5: when the region's accessible configuration space is
exhausted via direct sampling, the biographical record itself becomes
the substrate for re-reading under a different protocol. The substrate
for reinterpretation is \ensuremath{\mathcal{B}(R).}

\subsection*{6.5 Clarifications}

Three clarifications about the exhaustion concept are important to state
explicitly.

Exhaustion is regional, not global. In an infinite substrate, novel
configurations can always emerge in regions that have not yet exhausted.
The substrate as a whole has no exhaustion threshold; only bounded
regions do. The substrate continues producing novelty somewhere
indefinitely.

Exhaustion does not mean dissolution. A region in Phase III still hosts
existing configurations, including persisting subsystems whose internal
dynamics continue. What stops is the production of new compressible
structure. The region becomes a museum of already-sampled configurations
rather than a generator of new ones.

Exhaustion does not mean stasis. Existing subsystems can still compute
on their data; stressors continue to operate; some subsystems still
persist while others dissolve. The region is dynamic. It is just no
longer novel. Every state the region enters in Phase III is a state it
has been in before.

\subsection*{6.6 The Structural Problem}

At \ensuremath{\tau} \textgreater{} \ensuremath{\tau}\emph{(R), the region holds confirmed bits with no
remaining direct-read protocol that can extract novel information.
\ensuremath{\Sigma_{\mathrm{sampled}}(R,} \ensuremath{\tau}}) \ensuremath{\approx} \ensuremath{\Sigma (R);} the configurations have all been sampled;
further sampling produces only repetitions. The confirmed bits still
exist, however. They satisfy Axiom 4 (causal anchoring) --- they are
part of the causal record and cannot be overwritten. They satisfy Axiom
3 (signal conservation) --- the signals that produced them have been
consumed and locked into the confirmation record.

The region is therefore in a structurally constrained state: it holds
informationally rich confirmed bits, but the substrate's direct-read
protocol (sampling new configurations, evaluating their stability
against the coherence threshold) can extract no further novelty from
them. Axiom 3 forbids losing these bits; the signals that produced them
have been irrevocably consumed and the record stands. Axiom 4 forbids
overwriting them. The substrate cannot eliminate the exhausted region;
it must continue to host it. But direct sampling has nothing further to
extract.

The only structurally available move is to read the confirmed bits under
a different protocol --- one that treats the confirmed bits not as
configurations to evaluate for stability, but as something else. The
biographical record \ensuremath{\mathcal{B}(R)} is available as the substrate for this
reinterpretation. What that protocol is, and how the reinterpretation
produces fundamental fields, initial conditions, and the wave dynamics
of a simulated universe, is the central work of Paper 5.

\section*{VII. Hierarchical Lightspeeds}
\addcontentsline{toc}{section}{VII. Hierarchical Lightspeeds}

\subsection*{7.1 The \ensuremath{c^{(k)}} Formula Restated}

From Paper 3 \S\,6.1, the propagation speed at hierarchical level k in
absolute substrate units is

\ensuremath{*c^{(k)}} = \ensuremath{\prod _{j=1}^{k}} \ensuremath{(L_{p}^{(j)}} / \ensuremath{L_{p}^{(j-1)})} \ensuremath{\cdot}
\ensuremath{(T_{p}^{(j-1)}} / \ensuremath{T_{p}^{(j)})} \ensuremath{\cdot} c*

with \ensuremath{L_{p}^{(0)}} \ensuremath{\equiv} \ensuremath{\ell_{p}} and \ensuremath{T_{p}^{(0)}} \ensuremath{\equiv} \ensuremath{t_{p}.~\text{The}} propagation speed in
level-k units (level-k spacings per level-k flop) is unchanged across
levels: \ensuremath{1/\sqrt{2}.} Every level inherits the same FCC face-diagonal geometry
that produced \ensuremath{c_{\mathrm{TMS}}} = \ensuremath{1/\sqrt{2}} at the substrate level (Paper 1 \S\,6.5). The
propagation-speed argument depends only on the FCC connectivity, which
is scale-invariant per Paper 3 \S\,4.5.

What varies across levels is the absolute propagation speed expressed in
substrate units. The variation comes from the ratios \ensuremath{L_{p}(j)/L_{p}(j-1)}
and \ensuremath{T_{p}(j-1)/T_{p}(j)} at each level --- geometric quantities set by the
minimum program extent and minimum program period at level j.

\subsection*{7.2 Measurement Nodes at Level k Measure \ensuremath{c^{(k)}}}

From \S\,IV.3, every subsystem at hierarchical level k \ensuremath{\geq} 1 is a measurement
node. A measurement node at level k reads its surround through imprint,
computes via its operation kernel, and produces output at the
meta-tetrahedral confirmation. Its natural spatial unit is \ensuremath{L_{p}^{(k)};}
its natural temporal unit is \ensuremath{T_{p}^{(k)}.}

When a measurement node at level k measures the propagation speed of a
wave crossing its meta-interstice, it counts how many \ensuremath{L_{p}^{(k)}}
spacings the wave covers per \ensuremath{T_{p}^{(k)}} flops. By the lattice-relative
invariance, this ratio is always \ensuremath{1/\sqrt{2}} in level-k units. Translated to
substrate units, it equals \ensuremath{c^{(k)}.}

Every measurement node at level k therefore measures the propagation
speed of its level: \ensuremath{1/\sqrt{2}} in its own units, \ensuremath{c^{(k)}} in substrate units.
The substrate c is the propagation speed at the agent level (k = 0),
derived from the FCC face-diagonal geometry. Measurement nodes above
level 0 measure in their own scaled units and observe \ensuremath{c^{(k)};} they can
infer the substrate c by applying the \ensuremath{c^{(k)}} formula in reverse if the
\ensuremath{L_{p}^{(j)}} and \ensuremath{T_{p}^{(j)}} ratios at intermediate levels are known. The
substrate c is inferrable but not directly measurable by any measurement
node above level 0.

If our universe is itself a level-k structure, our measured speed of
light is \ensuremath{c^{(k)}} at our level rather than the substrate c.~Which level
we sit at, and what the substrate c is in our SI units, are questions
Paper 5's simulation work pins.

\subsection*{7.3 The \ensuremath{c^{(k)}} Sequence as Dispersion Spectrum}

The ratios \ensuremath{L_{p}(k)/L_{p}(k-1)} and \ensuremath{T_{p}(k)/T_{p}(k-1)} at each level are
independent geometric quantities. \ensuremath{L_{p}^{(k)}} is set by the minimum
spatial extent of a stabilizer-satisfying T-closed configuration at
level k; \ensuremath{T_{p}^{(k)}} is set by the minimum temporal period of such a
configuration. These ratios depend on different aspects of the level-k
combinatorics and are not constrained to scale together.

Consequence: \ensuremath{c^{(k)}} does not follow a monotonic trend across levels.
Level 1 might have \ensuremath{c^{(1)}} \textless{} c, if the spacing ratio rises
faster than the period ratio falls. Level 2 might have \ensuremath{c^{(2)}}
\textgreater{} \ensuremath{c^{(1)}} if the level-2 ratios produce a faster
propagation. Level 3 might be smaller again. The hierarchical
lightspeeds form a dispersion spectrum --- a sequence of distinct
absolute speeds, each set by local combinatorial facts about programs at
its level.

There is no general inequality \ensuremath{c^{(k+1)}} \textgreater{} \ensuremath{c^{(k)}} or
\ensuremath{c^{(k+1)}} \textless{} \ensuremath{c^{(k)}} without additional constraints. The sign
of \ensuremath{(c^{(k+1)}} \ensuremath{-} \ensuremath{c^{(k)})} depends on whether \ensuremath{(L_{p}(k+1)/L_{p}(k))} \ensuremath{\cdot}
\ensuremath{(T_{p}(k)/T_{p}(k+1))} is greater than or less than 1, which depends on the
level-(k+1) combinatorics independently of level-k.

\subsection*{7.4 Bounds on the Spectrum}

The dispersion is bounded above. The coupling-timescale constraint from
\S\,5.3 forces \ensuremath{T_{p}^{(k+1)}} \ensuremath{\geq} \ensuremath{\sqrt{2}} \ensuremath{\cdot} \ensuremath{(L_{p}(k+1)/L_{p}(k))} \ensuremath{\cdot} \ensuremath{T_{p}^{(k)}.}
Substituting into the \ensuremath{c^{(k+1)}} formula:

\emph{\ensuremath{c^{(k+1)}} \ensuremath{\leq} \ensuremath{(L_{p}}\textbf{\ensuremath{(k+1)/L_{p}}}(k)) \ensuremath{\cdot} \ensuremath{(T_{p}^{(k)}} / \ensuremath{[\sqrt{2}} \ensuremath{\cdot}
\ensuremath{(L_{p}}\textbf{\ensuremath{(k+1)/L_{p}}}(k)) \ensuremath{\cdot} \ensuremath{T_{p}^{(k)}])} \ensuremath{\cdot} \ensuremath{c^{(k)}} = \ensuremath{(1/\sqrt{2})} \ensuremath{\cdot}
\ensuremath{c^{(k)}}}

So \ensuremath{c^{(k+1)}} \ensuremath{\leq} \ensuremath{(1/\sqrt{2})} \ensuremath{\cdot} \ensuremath{c^{(k)}} is a derived upper bound. The \ensuremath{1/\sqrt{2}}
factor has a sharper geometric origin than the coupling-timescale
derivation above suggests: it is the FCC face-diagonal-to-edge ratio of
the cubic cell containing the lattice at every hierarchical level, made
explicit in Paper 5 \S\,7.2. By recursive scale-invariance and the
sublattice bifurcation of Paper 3 \S\,4.5, the same FCC
face-diagonal-to-edge ratio reappears at every level k \ensuremath{\geq} 0. The
\ensuremath{c^{(k+1)}} \ensuremath{\leq} \ensuremath{(1/\sqrt{2})} \ensuremath{c^{(k)}} bound is therefore not merely a constraint
derived from the coupling-timescale argument: it is a \emph{structural
law} of the recursive FCC geometry, with the coupling-timescale argument
here providing the dynamical mechanism by which the geometric law is
realized.

Lightspeed cannot grow across the hierarchy beyond this geometric
ceiling per level. There is no general lower bound --- \ensuremath{c^{(k+1)}} can be
arbitrarily small relative to \ensuremath{c^{(k)}} if the level-(k+1) ratios produce
a slow propagation. Lightspeeds can collapse downward through the
hierarchy; they cannot grow upward beyond the \ensuremath{(1/\sqrt{2})} factor per level.

The \ensuremath{c^{(k)}} sequence is therefore monotonically non-increasing in the
tight-coupling limit (where the upper bound is saturated), and freely
variable downward in looser coupling regimes. Higher hierarchical levels
propagate at most slightly slower than lower levels in the tightest
case, and potentially much slower in loose cases.

\subsection*{7.5 Combinatorial Origin of \ensuremath{L_{p}^{(k)}} and \ensuremath{T_{p}^{(k)}}}

The exact values of \ensuremath{L^{p(k)}} and \ensuremath{T^{p(k)}} are derived in Paper 5 \S\,6.1.3
via the sublattice-bifurcation recursive structure, with the empirical
verification of the bifurcation structure through Stages 5 and 6
documented in Paper 5 Appendix E. The substrate-level values are \ensuremath{L^{p(1)}}
= \ensuremath{3\sqrt{2}} \ensuremath{\ell_{p}} and \ensuremath{T^{p(1)}} = 6 \ensuremath{t_{p}} (Paper 5 \S\,6.1.3, \S\,6.2); higher-level
values follow the recursive ratio \ensuremath{L^{p(k)}} = \ensuremath{(3\sqrt{2})_{k}} \ensuremath{\ell_{p}} and \ensuremath{T^{p(k)}} = \ensuremath{6_{k}}
\ensuremath{t_{p}} in the saturation cascade. The combinatorial enumeration that
supports these values is documented in Paper 5 \S\,VI and Appendix B.

What Paper 4 establishes structurally: the ratios \ensuremath{L_{p}(k)/L_{p}(k-1)} and
\ensuremath{T_{p}(k)/T_{p}(k-1)} are finite and positive at every level (no level has
zero or infinite extent), they are independent of each other (set by
different aspects of stabilizer-satisfying configurations), and they are
bounded below by 1 + \ensuremath{\epsilon} for some \ensuremath{\epsilon} \textgreater{} 0 (the scale-separation
requirement of \S\,2.2 forces them away from 1). These structural facts are
sufficient to derive the dispersion-spectrum claim of \S\,7.3 and the
upper-bound claim of \S\,7.4 without pinning numerical values.

\subsection*{7.6 Empirical Consequences}

If our universe is itself a level-k structure for some k \ensuremath{\geq} 1 (with k
determined by Paper 5's simulation work), then our measured speed of
light is \ensuremath{c^{(k)}} in substrate units. We have no direct access to the
substrate c; our measurements are by construction in level-k units.
Three structural consequences follow.

Universal lightspeed within a level. All measurement nodes at our level
measure the same \ensuremath{c^{(k)}.} Lorentz invariance at our level is preserved:
the propagation-speed argument is scale-invariant per Paper 3 \S\,4.5, so
the lattice-relative invariance that gives \ensuremath{c^{(k)}} the same value to
all level-k measurement nodes is geometric and applies uniformly.

Different lightspeeds at different levels. Measurement nodes at level k'
\ensuremath{\neq} our k measure \ensuremath{c^{(k')},} which generically differs from our
c.~Cross-level interactions could produce apparent variations in
measured propagation speeds --- the substrate-level origins of
relativistic phenomena, flagged for Paper 5 development in \S\,VIII forward
references below.

The substrate c is inferrable but not directly measurable by any
measurement node above level 0. The substrate c is the propagation speed
at the agent level, derived from the FCC face-diagonal geometry; level-k
measurement nodes observe \ensuremath{c^{(k)}} and can compute the substrate c by
inverting the \ensuremath{c^{(k)}} formula given the intermediate \ensuremath{L_{p}^{(j)}} and
\ensuremath{T_{p}^{(j)}} ratios.

The hierarchical structure that produces \ensuremath{c^{(k)}} also produces
refinements to the GUP deformation parameter \ensuremath{\beta ,} developed in \S\,VIII.

\subsection*{7.7 Summary}

Hierarchical lightspeeds are not free parameters. Each \ensuremath{c^{(k)}} is
forced by the level's \ensuremath{L_{p}^{(k)}} and \ensuremath{T_{p}^{(k)}} ratios, which are
themselves forced by combinatorial enumeration of stabilizer-satisfying
configurations at that level. The sequence forms a dispersion spectrum
bounded above by \ensuremath{c^{(k+1)}} \ensuremath{\leq} \ensuremath{(1/\sqrt{2})} \ensuremath{\cdot} \ensuremath{c^{(k)}} and unbounded below.
Measurement nodes at each level measure their level's \ensuremath{c^{(k)}.} Our
measured c is \ensuremath{c^{(k)}} for some k determined by which hierarchical level
our universe occupies --- a question reserved for Paper 5's simulation
work.

\section*{VIII. Predictions and Forward References}
\addcontentsline{toc}{section}{VIII. Predictions and Forward References}

\subsection*{8.1 Refined GUP Prediction}

Paper 1 \S\,7.2 introduced \ensuremath{\beta} = f(d, \ensuremath{K_{\psi}):} the deformation parameter scales
with causal depth and polarization complexity. Paper 2 \S\,6.4 identified
\ensuremath{K_{\psi}} as the Kolmogorov complexity of the program running through the
bit's causal history. Paper 3 \S\,6.2 refined further: \ensuremath{\beta} = f(d,
\ensuremath{K_{\psi,\mathrm{local}},} \ensuremath{K_{\psi,\mathrm{surround}},} \ensuremath{K_{\psi,\mathrm{subsystem}}),} incorporating the
imprinting surround history and the encompassing subsystem context.
Paper 4 adds the final functional dependence:

\begin{equation*}
\beta = f(d, K_{\psi,\mathrm{local}}, K_{\psi,\mathrm{surround}}, K_{\psi,\mathrm{subsystem}}, k)
\end{equation*}

where k is the hierarchical level at which the bit's enclosing subsystem
operates. The function f is monotonically increasing in all five
arguments. A bit embedded in a level-k subsystem inherits the
computational complexity not only of its local region, surround, and
immediate subsystem, but also of the hierarchical depth at which that
subsystem participates.

The physical reading: deeper hierarchical embedding means more layers of
computational structure between the bit and the substrate level. Each
layer adds to the algorithmic complexity of the bit's causal record.
Bits in subsystems at higher k have larger substrate footprints --- more
lattice geometry and computational content is geometrically present in
their confirmation records. The closed form of f is reserved for Paper
5's simulation work, which can numerically probe the relationship
between hierarchical depth and effective \ensuremath{\beta} across the range of (d, \ensuremath{K_{\psi})}
values encountered in actual cascade instances. Paper 4 establishes the
functional dependence; Paper 5 will pin the form.

\subsection*{8.2 Cross-Level Propagation Effects}

\S\,VII derived the \ensuremath{c^{(k)}} sequence as a dispersion spectrum. The
structural and empirical consequences of \ensuremath{c^{(k)}} variation within a
level were treated. The consequences of \ensuremath{c^{(k)}} variation across levels
--- when subsystems at different hierarchical levels share substrate at
their interfaces --- are flagged here as forward reference.

When a \ensuremath{\text{level}-(k-1)} subsystem exists inside a level-k meta-subsystem
(downward Russian doll nesting from \S\,2.4), the two subsystems share
substrate at their interface. The \ensuremath{\text{level}-(k-1)} subsystem propagates
information at \ensuremath{c^{(k-1)};} the level-k subsystem propagates at \ensuremath{c^{(k)}.}
At the interface, two different propagation speeds meet on the same
substrate.

The substrate-level dynamics at this interface produce three structural
effects. Variable apparent propagation: a signal originating in the
\ensuremath{\text{level}-(k-1)} subsystem and observed at the level-k scale appears to
propagate at a level-mixed effective speed, with the mixing depending on
the boundary geometry and the relative \ensuremath{c^{(k)}} and \ensuremath{c^{(k-1)}}
magnitudes. Frame-relative timing: events that are simultaneous in
\ensuremath{\text{level}-(k-1)} units may not be simultaneous in level-k units, because the
temporal scale \ensuremath{T_{p}^{(k)}} differs from \ensuremath{T_{p}^{(k-1)};} measurement nodes
at different levels disagree about event ordering for events at
level-mixed locations. Embedded-medium dilation: a \ensuremath{\text{level}-(k-1)} subsystem
deeply embedded in a level-k meta-subsystem experiences propagation
through the meta-subsystem's substrate at the effective \ensuremath{c^{(k)}} of the
meta-subsystem's medium, producing an apparent slowing or speeding
relative to the bare \ensuremath{c^{(k-1)}.}

These three effects are the substrate-level origins of relativistic
phenomena. Their derivation from the \ensuremath{c^{(k)}} spectrum and the
cross-level interface dynamics requires the closed-form \ensuremath{c^{(k)}} values
and the field reinterpretation protocol of Paper 5, and is addressed
there. The QCA program (D'Ariano \& Perinotti 2014; Mlodinow \& Brun
2025) is a parallel direction deriving relativistic dynamics from
discrete substrates and provides useful comparison territory for that
work.

\subsection*{8.3 The Paper 5 Pivot}

Paper 4 \S\,VI established that any bounded substrate region traverses a
biography from formation through exhaustion. At \ensuremath{\tau} \textgreater{} \ensuremath{\tau *(R),}
the region's accessible configuration space is saturated; the
substrate's direct-read protocol can no longer extract novel
information; further sampling produces only repetitions of
already-sampled configurations.

The confirmed bits in R still exist. By Axiom 4 (causal anchoring), they
cannot be overwritten. By Axiom 3 (signal conservation), the signals
that produced them cannot be lost. The region holds informationally rich
confirmed bits with no remaining direct-read protocol that can extract
novelty from them.

The only structurally available move is to read the confirmed bits under
a different protocol. The substrate must reinterpret what the confirmed
bits represent --- treating them not as configurations to evaluate for
stability, but as something else. The biographical record itself, \ensuremath{\mathcal{B}(R),}
is available as the substrate for this reinterpretation.

Paper 5 develops the reinterpretation protocol. The work involves
identifying which protocols are structurally available given the axioms
(signal conservation, causal anchoring, generative economy); deriving
the specific reinterpretation that produces fundamental fields and
initial conditions; showing how the confirmed binary content becomes
wave configurations on those fields; deriving the spectrum of physical
laws that follow; and addressing the cross-level relativistic phenomena
flagged in \S\,8.2.

The pivot point is named here: at \ensuremath{\tau} \textgreater{} \ensuremath{\tau *(R),} the
substrate's read protocol transitions. What comes after the transition
is the work of Paper 5.

\subsection*{8.4 Other Forward References}

To Paper 5: the closed-form \ensuremath{c^{(k)}} values; the exact integer
\ensuremath{L_{p}^{(k)};} the C(4,2) = 6 combinatorial refinement; the full closed
form of f(d, \ensuremath{K_{\psi,\mathrm{local}},} \ensuremath{K_{\psi,\mathrm{surround}},} \ensuremath{K_{\psi,\mathrm{subsystem}},} k); the
cross-level relativistic derivation; the reinterpretation protocol; the
fundamental field identification; the initial conditions; the Big Bang
correspondence.

To Paper 6 (if scope migrates): the full quadratic Born rule including
complex amplitudes, phase, and interference; the
\ensuremath{K_{\psi}-to-\text{metric}-\text{uncertainty}} theorem.

To Volume 2: observer experience as the agent-side complement of
measurement-node structure; the relationship between region biography
\ensuremath{\mathcal{B}(R)} and subjective biography; the experiential reading of
self-modeling-of-self at k \ensuremath{\geq} 2; contact with the conscious-agents
framework of Hoffman, Prakash, Prentner (2023) and Integrated
Information Theory 4.0 (Albantakis et al.~2023); the
observer-in-cosmology program (Gorobey et al.~2025; Zhao, Harlow,
Usatyuk 2025) and the von-Neumann-algebra approach to observers as
bounded structures in closed universes.

\subsection*{8.5 Summary}

Paper 4 has established the macroscopic consequences of subsystem
dynamics. The recursive coarse-graining operator \ensuremath{G^{(k)}} lifts the
substrate's geometry to arbitrary hierarchical depth; the upward
recursion is bounded above by region extent; the downward recursion is
bounded below by \ensuremath{L_{p};} the two recursions are asymmetric, with each
level genuinely contributing independent configuration space. Subsystems
grow via annexation, imprinting, and incorporation, with three different
kinds of ceiling. Self-modeling is a structural property of every
subsystem at k \ensuremath{\geq} 1, and every such subsystem is a measurement node.
Multi-subsystem dynamics produce three derivable stressors ---
mortality, predation, starvation --- each with substrate-level mechanism
and level-(k+1) signature. Diversity arises both inter-region (different
stressor profiles) and intra-region (single stressor underdetermines its
survival set). Subsystem formation is overdetermined by Borel-Cantelli,
PGE preference, and flux-balance forcing. Region biography \ensuremath{\mathcal{B}(R)} traces a
bounded region from formation through exhaustion in three phases.
Novelty exhausts in finite flop time because \textbar \ensuremath{\Sigma (R)}\textbar{} is
finite. The \ensuremath{c^{(k)}} spectrum forms a dispersion bounded above by \ensuremath{(1/\sqrt{2})}
\ensuremath{\cdot} \ensuremath{c^{(k-1)}.} The GUP parameter \ensuremath{\beta} refines with hierarchical level.

No new axioms. No new free numerical parameters. Every quantity derived
from the five axioms and the Principle of Generative Economy as applied
recursively across the hierarchy.

The arc from Paper 1 through Paper 4 is the derivation, from five axioms
and one principle, of a substrate that produces hierarchical structures
of arbitrary depth, each populated by self-modeling measurement nodes
that learn, grow, stress one another, diversify, and ultimately exhaust
their local configuration spaces --- at which point the substrate
confronts the structural problem that Paper 5 solves.

\section*{Appendix A: Full Expansion of \ensuremath{\nu *(R)}}
\addcontentsline{toc}{section}{Appendix A: Full Expansion of \ensuremath{\nu *(R)}}

This appendix provides the full symbolic expansion of the exhaustion
threshold \ensuremath{\nu *(R)} introduced in \S\,6.3, with each dissolution-rate term
derived from the \S\,V stressor mechanisms.

\subsection*{A.1 Mortality Rate}

From \S\,5.5, mortality arises when a subsystem's coherence falls below its
threshold \textbar \ensuremath{\Pi_{\psi}}\textbar\ensuremath{*(L_{\mathrm{sub}}^{(k)})} due to accumulated
uncorrectable errors. The per-event uncorrectable error probability in a
tetrahedral confirmation unit at level k is approximately

\emph{\ensuremath{P_{\mathrm{unc}}^{(k)}} \ensuremath{\approx} C(4,2) \ensuremath{\cdot} \ensuremath{(p_{\mathrm{error}}^{(k)})^{2}} = 6 \ensuremath{\cdot} ((1 \ensuremath{-}
\textbar \ensuremath{\Pi_{\psi}}\textbar\ensuremath{^{(k)})/2)^{2}}}

where the factor C(4,2) = 6 is the worst-case upper bound on the number
of two-error configurations in a tetrahedral unit (Paper 3 \S\,3.3); the
exact refinement of this combinatorial factor is reserved for Paper 5.
Per level-k flop, the mortality rate is

\begin{equation*}
d_{\mathrm{mort}}^{(k)} \approx P_{\mathrm{unc}}^{(k)} \cdot N_{\mathrm{sub}}^{(k)}
\end{equation*}

where \ensuremath{N_{\mathrm{sub}}^{(k)}} is the population of level-k subsystems in the
region. Mortality scales linearly in subsystem count and quadratically
in the per-event error probability.

\subsection*{A.2 Predation Rate}

From \S\,5.6, predation requires a predator-prey encounter --- a boundary
interaction between an active subsystem and a vulnerable neighbor. The
encounter rate scales as the square of subsystem density per unit
volume:

\begin{equation*}
d_{\mathrm{pred}}^{(k)} \approx G_{\mathrm{pred}}^{(k)} \cdot (N_{\mathrm{sub}}^{(k)})^{2} / V_{R}^{(k)}
\end{equation*}

where \ensuremath{V_{R}^{(k)}} is the region's volume in level-k lattice units, and
\ensuremath{G_{\mathrm{pred}}^{(k)}} is a geometric proportionality constant capturing the
typical boundary-overlap geometry between adjacent level-k subsystems.
\ensuremath{G_{\mathrm{pred}}^{(k)}} depends on \ensuremath{L_{p}^{(k)}} and the local subsystem packing
density; its exact value at each hierarchical level is reserved for
Paper 5's simulation work.

\subsection*{A.3 Starvation Rate}

From \S\,5.7, starvation arises when the surround becomes uninformative.
Per level-k flop, the starvation rate is approximately

\begin{equation*}
d_{\mathrm{starv}}^{(k)} \approx G_{\mathrm{starv}}^{(k)} \cdot N_{\mathrm{sub}}^{(k)} \cdot (1 -
K(\psi_{\mathrm{surround}})/K_{\mathrm{max}})
\end{equation*}

where \ensuremath{G_{\mathrm{starv}}^{(k)}} is the geometric proportionality constant
capturing imprint geometry at level k (set by the imprint depth \ensuremath{d_{R}} and
the boundary configuration), \ensuremath{K(\psi_{\mathrm{surround}})} is the algorithmic
complexity of the surround within imprint depth, and \ensuremath{K_{\mathrm{max}}} =
\textbar \ensuremath{C_{\mathrm{surround}}}\textbar{} \ensuremath{\cdot} log 2 is the maximum possible surround
complexity. As surround complexity approaches \ensuremath{K_{\mathrm{max}},} starvation
pressure approaches zero; as surround complexity falls toward
log(volume), starvation pressure approaches \ensuremath{G_{\mathrm{starv}}^{(k)}} \ensuremath{\cdot}
\ensuremath{N_{\mathrm{sub}}^{(k)}.}

\subsection*{A.4 Formation Rate}

From \S\,5.1, formation rate at novelty production rate \ensuremath{\nu} is

\emph{\ensuremath{formation_{\mathrm{rate}}(R,} \ensuremath{\tau ,} \ensuremath{\nu )} \ensuremath{\approx} \ensuremath{\rho_{\mathrm{sample}}(R)} \ensuremath{\cdot} \ensuremath{(\nu} /
\textbar \ensuremath{\Sigma (R)}\textbar) \ensuremath{\cdot} \ensuremath{P_{\mathrm{coherent}}}}

where \ensuremath{\rho_{\mathrm{sample}}(R)} is the substrate sampling rate in R and \ensuremath{P_{\mathrm{coherent}}}
is the probability that a sampled stable configuration meets the
coherence threshold \textbar \ensuremath{\Pi_{\psi}}\textbar\ensuremath{*(L_{\mathrm{sub}}^{(k)})} at its
corresponding level.

\subsection*{A.5 Exhaustion Threshold}

Setting formation rate equal to total dissolution rate and solving for
\ensuremath{\nu :}

\emph{\ensuremath{\rho_{\mathrm{sample}}(R)} \ensuremath{\cdot} \ensuremath{(\nu}}(R) / \textbar \ensuremath{\Sigma (R)}\textbar) \ensuremath{\cdot} \ensuremath{P_{\mathrm{coherent}}} =
\ensuremath{\Sigma_{k}} \ensuremath{[d_{\mathrm{mort}}^{(k)}} + \ensuremath{d_{\mathrm{pred}}^{(k)}} + \ensuremath{d_{\mathrm{starv}}^{(k)}]*}

Therefore

\emph{\ensuremath{\nu}}(R) = \textbar \ensuremath{\Sigma (R)}\textbar{} \ensuremath{\cdot} \ensuremath{\Sigma_{k}} \ensuremath{[d_{\mathrm{mort}}^{(k)}} +
\ensuremath{d_{\mathrm{pred}}^{(k)}} + \ensuremath{d_{\mathrm{starv}}^{(k)}]} / \ensuremath{(\rho_{\mathrm{sample}}(R)} \ensuremath{\cdot} \ensuremath{P_{\mathrm{coherent}})*}

Substituting the rate expressions from A.1--A.3:

\emph{\ensuremath{\nu}}(R) = \textbar \ensuremath{\Sigma (R)}\textbar{} / \ensuremath{(\rho_{\mathrm{sample}}(R)} \ensuremath{\cdot} \ensuremath{P_{\mathrm{coherent}})} \ensuremath{\cdot}
\ensuremath{\Sigma_{k}} [6 \ensuremath{\cdot} ((1 \ensuremath{-} \textbar \ensuremath{\Pi_{\psi}}\textbar\ensuremath{^{(k)})/2)^{2}} \ensuremath{\cdot} \ensuremath{N_{\mathrm{sub}}^{(k)}} +
\ensuremath{G_{\mathrm{pred}}^{(k)}} \ensuremath{\cdot} \ensuremath{(N_{\mathrm{sub}}^{(k)})^{2}} / \ensuremath{V_{R}^{(k)}} + \ensuremath{G_{\mathrm{starv}}^{(k)}} \ensuremath{\cdot}
\ensuremath{N_{\mathrm{sub}}^{(k)}} \ensuremath{\cdot} (1 \ensuremath{-} \ensuremath{K(\psi_{\mathrm{surround}}(k))/K_{\mathrm{max}}(k))]*}

This is the full symbolic expression. Every term except the geometric
proportionality constants \ensuremath{G_{\mathrm{pred}}^{(k)},} \ensuremath{G_{\mathrm{starv}}^{(k)},} and the
C(4,2) = 6 upper bound is fully derived from inherited machinery. The
numerical pinning of these structural constants is reserved for Paper
5's simulation work, in keeping with the structural-vs-numerical
division of labor across the volume.

\section*{Acknowledgments}
\addcontentsline{toc}{section}{Acknowledgments}

The development of this volume was substantially aided by extended
collaboration with several large language models used as critical
sounding boards, pedagogical resources, formal-rigor checkers, and
external working memory across many sessions during 2024--2026.

\textbf{Gemini (Google DeepMind)} was the long-running primary
collaborator across the framework's conceptual development, providing
physics pedagogy, structural critique, and ongoing review of the ideas
that became Volume 1. The author's path into the technical literature,
and the working-out of the framework's core conceptual moves over a long
period preceding formalization, were carried by extended dialogue with
Gemini.

\textbf{Claude (Anthropic)} was the primary collaborator across the
formalization, derivation tightening, simulation work, citation
auditing, and editorial passes that brought the volume to its present
form, including the sublattice bifurcation analysis empirically verified
in Paper 5 Appendix E.

\textbf{ChatGPT (OpenAI)}, \textbf{DeepSeek}, \textbf{Grok (xAI)}, and
\textbf{Granite (IBM)} contributed additional structural feedback and
review at various points.

All theoretical commitments, the choice of axioms, the Principle of
Generative Economy, the sublattice bifurcation insight, and the
framework's overall architecture are the author's; the AI systems' role
was to teach, formalize, critique, verify, and document.

\section*{References}
\addcontentsline{toc}{section}{References}
\begin{reflist}
Albanese, L., Barra, A., Bianco, P., Durante, F., \& Pallara, D. (2024).
Hebbian Learning from First Principles. Journal of Mathematical Physics,
65, 113302. doi:10.1063/5.0197652.
\url{https://arxiv.org/abs/2401.07110}

Albantakis, L., Barbosa, L. S., Findlay, G., Grasso, M., Haun, A. M.,
Marshall, W., Mayner, W. G. P., Zaeemzadeh, A., Boly, M., Juel, B. E.,
Sasai, S., Fujii, K., David, I., Hendren, J., Lang, J. P., \& Tononi, G.
(2023). Integrated information theory (IIT) 4.0: Formulating the
properties of phenomenal existence in physical terms. PLoS Computational
Biology, 19(10), e1011465. doi:10.1371/journal.pcbi.1011465.
\url{https://arxiv.org/abs/2212.14787}

Almheiri, A., Dong, X., \& Harlow, D. (2015). Bulk locality and quantum
error correction in AdS/CFT. Journal of High Energy Physics, 2015(4),
163. doi:10.1007/JHEP04(2015)163. \url{https://arxiv.org/abs/1411.7041}

Bombelli, L., Lee, J., Meyer, D., \& Sorkin, R. D. (1987). Space-time as
a causal set. Physical Review Letters, 59(5), 521--524.
doi:10.1103/PhysRevLett.59.521.

Born, M. (1926). Zur Quantenmechanik der Stoßvorgänge. Zeitschrift für
Physik, 37(12), 863--867. doi:10.1007/BF01397477.

D'Ariano, G. M., \& Perinotti, P. (2014). Derivation of the Dirac
equation from principles of information processing. Physical Review A,
90(6), 062106. doi:10.1103/PhysRevA.90.062106.
\url{https://arxiv.org/abs/1306.1934}

Goldstein, H., Poole, C., \& Safko, J. (2002). Classical Mechanics (3rd
ed.). Addison-Wesley.

Gorard, J. (2020). Some Quantum Mechanical Properties of the Wolfram
Model. Complex Systems, 29(2), 537--598.
doi:10.25088/ComplexSystems.29.2.537.

Gorobey, N., Lukyanenko, A., \& Goltsev, A. V. (2025). Observer in
Quantum Cosmology. arXiv preprint.
\url{https://arxiv.org/abs/2506.06379}

Gottesman, D. (1997). Stabilizer Codes and Quantum Error Correction.
Ph.D.~Thesis, California Institute of Technology.
\url{https://arxiv.org/abs/quant-ph/9705052}

Hamilton, W. R. (1834). On a General Method in Dynamics. Philosophical
Transactions of the Royal Society, 124, 247--308.

Hebb, D. O. (1949). The Organization of Behavior: A Neuropsychological
Theory. Wiley.

Hoffman, D. D., Prakash, C., \& Prentner, R. (2023). Fusions of
Consciousness. Entropy, 25(1), 129. doi:10.3390/e25010129.

Hossenfelder, S. (2013). Minimal Length Scale Scenarios for Quantum
Gravity. Living Reviews in Relativity, 16(1), 2.
doi:10.12942/lrr-2013-2.

Jaynes, E. T. (1957). Information Theory and Statistical Mechanics.
Physical Review, 106(4), 620--630. doi:10.1103/PhysRev.106.620.

Kauffman, S. (1993). The Origins of Order. Oxford University Press.

Kolmogorov, A. N. (1965). Three approaches to the quantitative
definition of information. Problems of Information Transmission, 1(1),
1--7.

Mlodinow, L., \& Brun, T. (2025). Quantum Electrodynamics from Quantum
Cellular Automata, and the Tension Between Symmetry, Locality, and
Positive Energy. Entropy, 27(5), 492. doi:10.3390/e27050492.
\url{https://arxiv.org/abs/2503.05998}

Pastawski, F., Yoshida, B., Harlow, D., \& Preskill, J. (2015).
Holographic quantum error-correcting codes: toy models for the
bulk/boundary correspondence. Journal of High Energy Physics, 2015(6),
149. doi:10.1007/JHEP06(2015)149. \url{https://arxiv.org/abs/1503.06237}

Rissanen, J. (1978). Modeling by shortest data description. Automatica,
14(5), 465--471. doi:10.1016/0005-1098(78)90005-5.

Rosas, F. E., Mediano, P. A. M., Luppi, A. I., Varley, T. F., Lizier, J.
T., Stramaglia, S., Jensen, H. J., \& Marinazzo, D. (2024). Software in
the natural world: A computational approach to hierarchical emergence.
arXiv preprint. \url{https://arxiv.org/abs/2402.09090}

Shannon, C. E. (1948). A mathematical theory of communication. Bell
System Technical Journal, 27(3), 379--423.
doi:10.1002/j.1538-7305.1948.tb01338.x.

Surya, S. (2019). The causal set approach to quantum gravity. Living
Reviews in Relativity, 22(1), 5. doi:10.1007/s41114-019-0023-1.

Switzer, A. (2026a). The Triadic Measurement Substrate (TMS), Volume 1,
Paper 1: Emergent 3D Information Topology and the Binary Alphabet via
[[4,2,2]] Quantum Error Detection Structural Correspondence.

Switzer, A. (2026b). The Triadic Measurement Substrate (TMS), Volume 1,
Paper 2: Computational Dynamics of the Binary Alphabet.

Switzer, A. (2026c). The Triadic Measurement Substrate (TMS), Volume 1,
Paper 3: Computation Emerges.

Takayanagi, T. (2025). Emergent Holographic Spacetime from Quantum
Information. Physical Review Letters, 134, 240001.
doi:10.1103/pg4r-fy8n.

't Hooft, G. (2016). The Cellular Automaton Interpretation of Quantum
Mechanics. Fundamental Theories of Physics 185. Springer.
doi:10.1007/978-3-319-41285-6.

Wheeler, J. A. (1990). Information, physics, quantum: The search for
links. In W. H. Zurek (Ed.), Complexity, Entropy and the Physics of
Information (pp.~3--28). Addison-Wesley.

Wolfram, S. (2002). A New Kind of Science. Wolfram Media.

Wolfram, S. (2020). A Class of Models with the Potential to Represent
Fundamental Physics. Complex Systems, 29(2), 107--536.
doi:10.25088/ComplexSystems.29.2.107.

Zhao, E. Y., Harlow, D., \& Usatyuk, M. (2025). Observers in a closed
universe. arXiv preprint.
\end{reflist}

\clearpage


\end{document}
